\documentclass[12pt]{article}

% Packages for mathematics
\usepackage{amsmath}
\usepackage{amssymb}
\usepackage{amsthm}
\usepackage{enumitem}
\usepackage{setspace}
\usepackage{url} 
\usepackage{caption}
\usepackage{lipsum} % For placeholder text
\usepackage{booktabs} % For better table formatting
\usepackage{booktabs}
\usepackage{graphicx}
\usepackage{subcaption}
\usepackage{float}
\usepackage{authblk}
\usepackage{makecell}

% Redefine \toprule, \midrule, and \bottomrule with custom thickness
\newcommand{\CellWithForceBreak}[2][c]{
\begin{tabular}[#1]{@{}c@{}}#2\end{tabular}}
\renewcommand{\toprule}{\specialrule{0.1mm}{0em}{0em}}  % Adjust thickness here
\renewcommand{\midrule}{\specialrule{0.1mm}{0em}{0em}}  % Adjust thickness here
\renewcommand{\bottomrule}{\specialrule{0.1mm}{0em}{0em}}  % Adjust thickness here

% Packages for graphics and tables
\usepackage{graphicx}
\usepackage{booktabs}

% Bibliography
\usepackage{natbib}
\bibliographystyle{apalike}

% Define a new environment for quotes with automatic numbering
\newcounter{prompt}
\newenvironment{myprompt}[1][]{%
    \vspace{2em} % Add space before the caption
    \refstepcounter{prompt} % Increment the prompt counter
    \noindent\textbf{Prompt \theprompt:} \textit{#1} \\ \vspace{-1.5em} % Add the caption with prompt number and custom text
    \begin{quote} \small % Start the quote environment with smaller text
}{%
    \end{quote} % End the quote environment
}
\newcommand{\mydot}{\, \cdot \,}


\title{paneltime: Mathematical documentation}

\author{
  Espen Sirnes\thanks{Tromsø University Business School, University of Tromsø, The Arctic University, Tromsø, Norway.\\
\textbf{JEL Codes:}  \\
\textbf{Keywords:} 
  }
 

}


\date{August 2024}
\begin{document}
\onehalfspacing

\maketitle

In the following, we present the mathematical foundations underlying the code in the \texttt{paneltime} package. \texttt{paneltime} estimates panel data models with a dynamic ARIMA-type mean equation and a GARCH-type conditional variance equation, jointly with random or fixed effects, by maximum likelihood. We begin with fundamental definitions, including the normal distribution. While the reader is likely familiar with this, building from first principles clarifies all relevant variables and assumptions.

\section{Fundamentals}\label{the-ll}

We assume that the standardized residuals $\varepsilon_{i,t}/\sigma_{i,t}$ are normally distributed, so that the conditional density of $\varepsilon_{i,t}$ is
\begin{equation}
f\bigl(\varepsilon_{i,t}, \sigma_{i,t}\bigr)
= \frac{1}{\sqrt{2\pi \,\sigma_{i,t}^{2}}}
  \exp\!\Bigl(- \tfrac{1}{2}
      \frac{\varepsilon_{i,t}^{2}}{\sigma_{i,t}^{2}}\Bigr),
\label{eq:normal_pdf}
\end{equation}
The term $\varepsilon_{i,t}$ denotes the innovation after differencing (if requested), the ARMA filter, and the group/time effect transformation. These residuals correspond to group $i$, out of $N$ total groups, at sequence position $t$, where $t$ ranges up to a maximum of $T$. The dataset may be unbalanced, meaning some groups can have fewer than $T$ observations. The function $\exp(\cdot)$ denotes the exponential function, and $\sigma_{i,t}^{2}$ is the conditional variance from a GARCH process \citep{engle1982autoregressive,bollerslev1986generalized}, specific to group $i$ at time $t$.

For the default Gaussian likelihood, the implementation uses $\varepsilon_{i,t}^{2}+10^{-8}$ in place of $\varepsilon_{i,t}^{2}$ in the quadratic likelihood term, and clips conditional variances to $[10^{-30},10^{30}]$ for numerical stability. These safeguards are implementation details; Equation~\ref{eq:normal_pdf} gives the statistical model.

In the remainder of the documentation, we will generally refer to the matrix representation of variables. That is, instead of working with scalar forms $\varepsilon_{i,t}$ and $\sigma_{i,t}^{2}$, we use their matrix counterparts $\mathbf{\varepsilon}$ and $\boldsymbol{\sigma}^{2}$. For instance, $\mathbf{\varepsilon}$ is an $N \times T$ matrix, where the entry in row $i$ and column $t$ corresponds to $\varepsilon_{i,t}$.

In cases where a group has fewer than $T$ observations, its observed values are placed at the beginning of the corresponding row, and the remaining entries are padded with zeros. Therefore, the matrix column index does not correspond to a common calendar date across groups. Any correction for RE or FE along the time dimension must account for this, using the actual date structure rather than assuming aligned positions across rows.

\subsection{Observed Residuals}

The regression residuals on the current panel arrays are defined as:
\begin{equation}
\mathbf{u}
= \mathbf{y} - \mathbf{X} \boldsymbol{\beta},
\label{eq:udef}
\end{equation}
where $\mathbf{X}$ is a three-dimensional array of shape $N \times T \times K$, containing the regressors (including the intercept column when present) across $N$ groups and $T$ within-group sequence positions. $\boldsymbol{\beta}$ is a column vector of regression coefficients, and $\mathbf{y}$ denotes the dependent variable. For $d>0$, these arrays have already undergone the differencing preprocessing described above; the first $d$ entries are zeroed and remain in the inclusion mask, while trailing padding is excluded.

The likelihood starts from the regression residual $\mathbf{u}=\mathbf{y}-\mathbf{X}\boldsymbol{\beta}$. When $d>0$, the panel arrays for $\mathbf{y}$ and the non-intercept regressors are differenced before this residual is formed; the intercept column is restored after differencing and the first $d$ positions of each group are set to zero. The ARMA filter is then applied to $\mathbf{u}$, the group/time effect corrections are applied to the resulting ARMA innovations, the GARCH recursion is applied to the effect-corrected innovations, and likelihood contributions are summed over observed cells.

The parameters of these transformations are estimated jointly by maximum likelihood. For unbalanced panels, rows are sorted within group and padded with zeros at the end; ARMA/GARCH lags follow that within-group observation order. Group/time effect means, in contrast, use actual group membership and the calendar-date map.


\subsection{ARIMA Correction}

Serial dependence is removed via a Box--Jenkins ARIMA($p,d,q$) filter \citep{box1970time}:
\begin{equation}
\tilde{\boldsymbol{\varepsilon}}
= \mathbf{MA}^{-1} \cdot \mathbf{AR} \cdot \mathbf{u},
\label{eq:arima_reduced}
\end{equation}
where $\mathbf{u}$ is formed from the panel inputs after differencing (if $d>0$), and $\tilde{\boldsymbol{\varepsilon}}$ is the resulting ARMA innovation matrix before the group/time effect corrections. The matrices $\mathbf{MA}$ and $\mathbf{AR}$ are lower-triangular lag operators applied separately to each group row. Differencing is a preprocessing step; it is not part of this ARMA equation.

\subsection{Random/Fixed Effects}

The group and/or time effects are removed from the ARMA innovations by centered quasi-demeaning corrections:

\begin{equation}
\varepsilon_{i,t}
= \tilde{\varepsilon}_{i,t}
 + \theta_{g,i}\left(\bar{\tilde{\varepsilon}}-\bar{\tilde{\varepsilon}}_{g,i}\right)
 + \theta_{t,t}\left(\bar{\tilde{\varepsilon}}-\bar{\tilde{\varepsilon}}_{t}\right),
\label{eq:rRE_def}
\end{equation}
where $\bar{\tilde{\varepsilon}}$ is the grand mean over included observations, $\bar{\tilde{\varepsilon}}_{g,i}$ is the mean for group $i$, and $\bar{\tilde{\varepsilon}}_{t}$ is the mean for actual calendar date $t$. A term is omitted when that effect is \texttt{none}. For \texttt{fixed} effects its weight is one; for \texttt{random} effects the weight $\theta$ is estimated from variance components, as described in Section~\ref{sec:effects-implementation}. In the two-way case, group and time corrections are each computed from the same $\tilde{\varepsilon}$ and added; they are not applied sequentially. Thus $\boldsymbol{\varepsilon}$ is a transformed innovation, not an innovation plus separately estimated $v_i$ and $w_t$ coefficients. The RE transformation approaches FE as its estimated effect variance grows relative to the idiosyncratic variance.

The effects are removed after the ARMA filter for two reasons. First, since all groups share the same filter, the effects keep their structure in $\tilde{\boldsymbol{\varepsilon}}$: if $u_{i,t}=v_i+w_t+\eta_{i,t}$, the filtered group effect is, after the first $\max(p,q)$ positions, the constant $v_i\left(1-\sum_j\rho_j\right)/\left(1+\sum_j\lambda_j\right)$ within each group, and in a balanced panel the filtered time effect is common to all groups at each date. Second, the GARCH recursion then operates on innovations from which the effects have been removed, so the effects are not interpreted as volatility. The correction is not exact during the first $\max(p,q)$ positions of each group, and the filtered time effect is serially correlated (for AR(1), $w_t-\rho w_{t-1}$), which scalar quasi-demeaning does not remove.

The transform is thus exact when the effects enter the innovation equation, $\mathbf{u}=\mathbf{AR}^{-1}\mathbf{MA}\,(v_i+w_t+\boldsymbol{\eta})$, since the filter then returns $v_i+w_t+\eta_{i,t}$ directly. When the effects are added to $\mathbf{u}$ itself, as above, it is an approximation. In simulations with ARMA(2,2)--GARCH(2,2) residuals, unit-variance group and time effects, $N=20$ and $T=100$, the ARMA estimates were essentially the same whether the effects were added to $\mathbf{u}$ or entered the innovation equation (bias about $0.04$ in both cases), so the approximation error is small relative to the ordinary finite-sample error.

\subsection{GARCH Correction}

Heteroscedasticity is then corrected by a GARCH(k,m) transformation \citep{bollerslev1986generalized}:
\begin{equation}
\mathbf{\sigma}^{2}
= \mathbf{GA}_{\sigma}^{-1}\,\mathbf{G}
  \;+\;
  \mathbf{GA}_{\sigma}^{-1}\,\mathbf{AR}_{\sigma}\,h(\mathbf{\varepsilon},z)
\label{eq:garch_reduced}
\end{equation}
where $\mathbf{\sigma}^{2}$ and $h(\mathbf{\varepsilon},z)$ are $N \times T$ arrays, and $\mathbf{GA}_{\sigma}$ and $\mathbf{AR}_{\sigma}$ are $T \times T$ GARCH and ARCH lag operators. The variance-regressor design $\mathbf{W}$ includes a constant and any user-supplied variance predictors, and $\mathbf{G}=\mathbf{W}\boldsymbol{\omega}$ except at the first observation of each group, where the implementation substitutes the initial variance. That initial variance is fixed from the starting-value calculation unless \texttt{include\_initvar=True}, in which case it is an estimated parameter. Presample innovations and variances are initialized at zero. For the default Gaussian model, $h(\varepsilon)=\varepsilon^2+10^{-8}$; custom h-functions and EGARCH use their configured transformations. The likelihood uses the resulting standardized innovations, subject to the numerical variance bounds described above.

\section{The ARIMA Process}\label{the-arima-process}

As described above, $\mathbf{u}$ contains regression residuals after differencing (if $d>0$). These residuals are linked to the ARMA innovations $\tilde{\boldsymbol{\varepsilon}}$ through:

\begin{equation}
\left( \mathbf{I} - \sum_{j = 1}^{p} \rho_{j} \cdot \mathbf{L}^{j} \right) \cdot \mathbf{u}
= \left( \mathbf{I} + \sum_{j = 1}^{q} \lambda_{j} \cdot \mathbf{L}^{j} \right) \cdot \tilde{\boldsymbol{\varepsilon}},
\label{eq:7}
\end{equation}

where $\mathbf{L}$ is the lag operator (or lag matrix). The moving average component is defined as:

\begin{equation}
\mathbf{MA} = \mathbf{I} + \sum_{j = 1}^{q} \lambda_{j} \cdot \mathbf{L}^{j},
\label{eq:8}
\end{equation}

and the autoregressive component is defined as:

\begin{equation}
\mathbf{AR} = \mathbf{I} - \sum_{j = 1}^{p} \rho_{j} \cdot \mathbf{L}^{j}.
\label{eq:9}
\end{equation}

We can then rewrite Equation~\ref{eq:7} as:

\begin{equation}
\tilde{\boldsymbol{\varepsilon}}
= \mathbf{MA}^{-1} \cdot \mathbf{AR} \cdot \mathbf{u},
\label{eq:arima_matrix}
\end{equation}

which expresses the innovations $\tilde{\boldsymbol{\varepsilon}}$ as a transformation of the already-differenced residuals $\mathbf{u}$.

\subsection*{Gradients}
The derivatives in this section are those of the ARMA innovations $\tilde{\boldsymbol{\varepsilon}}$, before the group/time effect corrections. The derivatives of the effect-corrected innovations $\boldsymbol{\varepsilon}$ used in the likelihood follow from these through the effect transform, as given in Equations~\ref{eq:re-two-way-first} and~\ref{eq:re-two-way-second}. With fixed or no effects, the transform is linear with constant weights, so the derivatives of $\boldsymbol{\varepsilon}$ are the same transform applied to the derivatives of $\tilde{\boldsymbol{\varepsilon}}$.

When writing down the gradients, it is convenient to define the following function
\begin{equation}
\mathbf{D}(n)
= - \mathbf{MA}^{-1} \cdot [\mathbf{L}^{1}, \cdots ,\mathbf{L}^{n}]
\label{eq:MA_func}
\end{equation}
The partial derivatives of $\tilde{\boldsymbol{\varepsilon}}$ with respect to the model parameters are:

\begin{equation}
\frac{\partial \tilde{\boldsymbol{\varepsilon}}}{\partial \boldsymbol{\rho}}
= \mathbf{D}(p) \, \mathbf{u}
\label{eq:11}
\end{equation}

\begin{equation}
\frac{\partial \tilde{\boldsymbol{\varepsilon}}}{\partial \boldsymbol{\lambda}}
= \mathbf{D}(q) \, \tilde{\boldsymbol{\varepsilon}}
\label{eq:12}
\end{equation}
Note that the derivative with respect to $\boldsymbol{\lambda}$ is taken with respect to $\tilde{\boldsymbol{\varepsilon}}$ itself rather than $\mathbf{u}$: since $\mathbf{MA}$ is inverted in Equation~\ref{eq:arima_reduced}, differentiating $\mathbf{MA}^{-1}$ with respect to $\lambda_j$ reintroduces a factor of $\mathbf{MA}^{-1}$, i.e.\ $\partial \mathbf{MA}^{-1}/\partial \lambda_j = -\mathbf{MA}^{-1}\mathbf{L}^{j}\mathbf{MA}^{-1}$, so that $\partial \tilde{\boldsymbol{\varepsilon}}/\partial \lambda_j = -\mathbf{MA}^{-1}\mathbf{L}^{j}\tilde{\boldsymbol{\varepsilon}}$.


\begin{equation}
\frac{\partial \tilde{\boldsymbol{\varepsilon}}}{\partial \boldsymbol{\beta}}
= - \mathbf{MA}^{-1} \cdot \mathbf{AR} \cdot \mathbf{X}
\label{eq:13}
\end{equation}

\subsection{The Hessians}

The Hessians can now be written as

\begin{equation}
\frac{\partial^{2} \tilde{\boldsymbol{\varepsilon}}}{\partial \boldsymbol{\rho}' \partial \boldsymbol{\rho}}=
\frac{\partial^{2} \tilde{\boldsymbol{\varepsilon}}}{\partial \boldsymbol{\beta}' \partial \boldsymbol{\beta}}
= 0
\label{eq:14}
\end{equation}
The mixed $\boldsymbol{\rho}$--$\boldsymbol{\beta}$ derivative is generally nonzero; it is given by Equation~\ref{eq:17}.

\begin{equation}
\frac{\partial^{2} \tilde{\boldsymbol{\varepsilon}}}{\partial \boldsymbol{\lambda}' \partial \boldsymbol{\lambda}}
= 2  \mathbf{D}(q)' \cdot \frac{\partial \tilde{\boldsymbol{\varepsilon}}}{\partial \boldsymbol{\lambda}}
\label{eq:15}
\end{equation}

\begin{equation}
\frac{\partial^{2} \tilde{\boldsymbol{\varepsilon}}}{\partial \boldsymbol{\lambda}' \partial \boldsymbol{\rho}}
= \mathbf{D}(q)' \cdot \frac{\partial \tilde{\boldsymbol{\varepsilon}}}{\partial \boldsymbol{\rho}}
\label{eq:16}
\end{equation}
\begin{equation}
\frac{\partial^{2} \tilde{\boldsymbol{\varepsilon}}}{\partial \boldsymbol{\lambda}' \partial \boldsymbol{\beta}}
= \mathbf{D}(q)' \cdot \frac{\partial \tilde{\boldsymbol{\varepsilon}}}{\partial \boldsymbol{\beta}}
\label{eq:18}
\end{equation}

\begin{equation}
\frac{\partial^{2} \tilde{\boldsymbol{\varepsilon}}}{\partial \boldsymbol{\rho}' \partial \boldsymbol{\beta}}
=- \mathbf{D}(p)' \cdot \mathbf{X},
\label{eq:17}
\end{equation}


\section{The GARCH process}\label{the-garch-process}

The underlying objective of the GARCH transformation is to model serially dependent conditional heteroskedasticity \citep{engle1982autoregressive,bollerslev1986generalized}. What is required is therefore a kernel that captures this serial dependence. Hence, the objective of GARCH is to obtain the serially correlated conditional standard deviations $\mathbf{\sigma}$ that deflate the stationarized ARIMA residuals $\mathbf{\varepsilon}$ into white noise $\mathbf{\varepsilon}/\mathbf{\sigma}$.

This is, in a sense, the opposite of the objective of the ARIMA transformation, which tries to \textit{remove} serial correlation in the conditional mean.


We shall assume a general GARCH($k,m$) process of the form

\begin{equation}
\mathbf{\sigma}^{2} = \mathbf{G} + \sum_{j = 1}^{m}{\psi_{j}\mathbf{\mathbf{L}}^{j}h\left( \mathbf{\varepsilon},z \right)} + \sum_{j = 1}^{k}{\gamma_{j} \cdot \mathbf{\mathbf{L}}^{j}\mathbf{\sigma}^{2}}
\label{eq:19}
\end{equation}
For the ordinary GARCH model $h( \mathbf{\varepsilon} ) = \mathbf{\varepsilon}^{2}$, and for an EGARCH-type model \citep{nelson1991conditional} $h( \mathbf{\varepsilon} ) = \ln\left( \mathbf{\varepsilon}^{2} \right)$. In the package, the user can specify the function $h( \cdot )$, which permits substantial freedom in the specification of the variance model. The user also has the option to let $h(\cdot)$ depend on a shape parameter $z$, which is then estimated jointly with the other parameters by maximizing the log likelihood.

The default \emph{h}-function is the usual GARCH model:

\begin{equation}
h( \mathbf{\varepsilon} ) = \mathbf{\varepsilon}^{2} 
\label{eq:20}
\end{equation}
If a different $h( \cdot )$ is supplied, its first and second derivative with respect to $\mathbf{\varepsilon}$ and, if used, $z$, must be defined also.

$\mathbf{W}$ contains a variance constant and any supplied variance predictors. The initial variance replaces, rather than adds to, the variance-regressor contribution at the first included observation of each group. Writing that initial position as $t=1$ within each row,
\begin{equation}
G_{i,t}=\begin{cases}
\sigma_{0}^{2}, & t=1,\\
W_{i,t}\boldsymbol{\omega}, & t>1.
\end{cases}
\label{eq:24b}
\end{equation}
The initial variance is fixed at its starting value by default; \texttt{include\_}\allowbreak\texttt{initvar=True} adds it as an estimated parameter. In the recursions, presample innovations and variances are initialized at zero.

The equation above can be reformulated as


\begin{equation}
\left( 1 - \sum_{j = 1}^{k}{\gamma_{j} \cdot}\mathbf{L}^{j} \right)\mathbf{\sigma}^{2} = \mathbf{G} + \sum_{j = 1}^{m}{\psi_{j}\mathbf{\mathbf{L}}^{j}h\left( \mathbf{\varepsilon},z \right)}
\label{eq:21}
\end{equation}

Let

\begin{equation}
\mathbf{GA}_{\sigma} =  I - \sum_{j = 1}^{k}{\gamma_{j} \cdot \mathbf{L}^{j}}
\label{eq:22}
\end{equation}

\begin{equation}
\mathbf{AR}_{\sigma} = \sum_{j = 1}^{m}{\psi_{j} \cdot \mathbf{L}^{j}}
\label{eq:23}
\end{equation}

Then the closed form solution of the GARCH process in matrix notation is, similar to the ARIMA model
\begin{equation}
 \mathbf{\sigma}^{2}  = \mathbf{GA}_{\sigma}^{- 1}\cdot \left(\mathbf{G} +\mathbf{AR}_{\sigma} \cdot h( \mathbf{\varepsilon},z ) \right)
\label{eq:24}
\end{equation}
Equations~\ref{eq:19}--\ref{eq:24} use the variance scale and describe ordinary GARCH. For \texttt{vol=EGARCH}, the implementation instead recurses on the log-variance state $\boldsymbol{\ell}=\log\boldsymbol{\sigma}^{2}$:
\begin{equation}
\boldsymbol{\ell}
=\mathbf{G}+\sum_{j=1}^{m}\psi_j\mathbf{L}^j\log(\boldsymbol{\varepsilon}^{2}+10^{-8})
+\sum_{j=1}^{k}\gamma_j\mathbf{L}^j\boldsymbol{\ell},
\qquad \boldsymbol{\sigma}^{2}=\exp(\boldsymbol{\ell}).
\label{eq:egarch-log-state}
\end{equation}
The Gaussian likelihood is evaluated using $\boldsymbol{\sigma}^{2}=\exp(\boldsymbol{\ell})$.

\subsection{A note on terminology}
The labels ARCH and GARCH follow the established terminology of \citet{engle1982autoregressive} and \citet{bollerslev1986generalized}. However, the word ``autoregressive'' in the name ARCH should not be read as saying that the $h(\boldsymbol{\varepsilon},z)$ term in Equation~\ref{eq:19} is itself an autoregressive feedback process. The $\boldsymbol{\psi}$ lag polynomial applies transformed past innovations to the variance equation. Since $h$ is a function of innovations rather than lagged conditional variances, this term describes the short-run response to shocks and is analogous in role to the moving-average side of an ARMA model; it is not literally an MA process, because $h$ may be nonlinear.

By contrast, the $\boldsymbol{\gamma}$ lag polynomial acts on past conditional variances and feeds them back into the current variance. This is the autoregressive, persistence-producing part of the recursion and is analogous in role to the AR side of an ARMA model. The $\boldsymbol{\psi}$ term is analogous to the MA-side shock response, while remaining a distributed lag of the generally nonlinear transform $h(\boldsymbol{\varepsilon},z)$ rather than a literal MA process. We retain the conventional ARCH and GARCH names, while using $\mathbf{AR}_\sigma$ for the innovation-response operator and $\mathbf{GA}_\sigma$ for the lagged-variance feedback operator in Equation~\ref{eq:garch_reduced}.


\subsection{Gradients}
When writing down the gradients, it is convenient to define the following function, analogous to $\mathbf{D}(n)$ in Equation~\ref{eq:MA_func}:
\begin{equation}
\mathbf{B}(n)
= \mathbf{GA}_{\sigma}^{-1} \cdot [\mathbf{L}^{1},  \cdots , \mathbf{L}^{n}]
\label{eq:GA_func}
\end{equation}
The first derivatives are then:
\begin{equation}
\frac{\partial \mathbf{\sigma}^{2} }{\partial\mathbf{\psi}} =\mathbf{B}(m) \cdot h( \mathbf{\varepsilon},z )
\label{eq:26}
\end{equation}
\begin{equation}
\frac{\partial \mathbf{\sigma}^{2} }{\partial\boldsymbol{\omega}} =  \mathbf{GA}_{\sigma}^{-1} \cdot
\mathbf{W}^{*}
\label{eq:27}
\end{equation}
where $\mathbf{W}^{*}$ is $\mathbf{W}$ with its first included entry in each group row set to zero, since that entry is replaced by the initial variance.
\begin{equation}
\frac{\partial \mathbf{\sigma}^{2} }{\partial\sigma_0^{2}} =\mathbf{GA}_{\sigma}^{-1} \cdot
\mathbf{e}
\label{eq:28}
\end{equation}
Here $\mathbf{e}$ has ones at the first included position in each group row and zeros elsewhere; this derivative is used only when the initial variance is estimated with \texttt{include\_initvar=True}.
\begin{equation}
\frac{\partial \mathbf{\sigma}^{2} }{\partial z} =\mathbf{GA}_{\sigma}^{-1} \mathbf{AR}_{\sigma} \cdot \frac{\partial h}{\partial z}
\label{eq:29}
\end{equation}
\begin{equation}
\frac{\partial \mathbf{\sigma}^{2} }{\partial\mathbf{\varepsilon}} = \mathbf{GA}_{\sigma}^{- 1}\cdot \mathbf{AR}_{\sigma} \cdot \frac{\partial h}{\partial \mathbf{\varepsilon}}
\label{eq:32}
\end{equation}
For $\boldsymbol{\gamma}$, which enters through the inverted operator $\mathbf{GA}_{\sigma}$ itself (the GARCH counterpart of $\boldsymbol{\lambda}$ in Equation~\ref{eq:12}), differentiating $\mathbf{GA}_{\sigma}^{-1}$ reintroduces a factor of $\mathbf{GA}_{\sigma}^{-1}$, giving
\begin{equation}
\frac{\partial \mathbf{\sigma}^{2} }{\partial\mathbf{\gamma}} = \mathbf{B}(k)\cdot\boldsymbol{\sigma}^{2}
\label{eq:25}
\end{equation}
(not $\mathbf{B}(k)\cdot(\mathbf{G}+\mathbf{AR}_{\sigma}\,h(\boldsymbol{\varepsilon},z))$: since $\boldsymbol{\sigma}^{2}=\mathbf{GA}_{\sigma}^{-1}(\mathbf{G}+\mathbf{AR}_{\sigma}h)$ already contains one factor of $\mathbf{GA}_{\sigma}^{-1}$, and this does not cancel against the $\mathbf{GA}_{\sigma}^{-1}$ inside $\mathbf{B}(k)$).

\subsection{Hessians}
The corresponding second derivatives, used in the Hessian, are
\begin{equation}
\frac{{\partial^{2} \mathbf{\sigma}^{2} }}{\partial\mathbf{\psi'}\partial\mathbf{\psi}} = 0
\label{eq:38}
\end{equation}
\begin{equation}
\frac{{\partial^{2} \mathbf{\sigma}^{2} }}{\partial \mathbf{\psi'} \, \partial z}
= \mathbf{B}(m) \, h_{z}\left( \mathbf{\varepsilon}, z \right)
\label{eq:39}
\end{equation}
\begin{equation}
\frac{{\partial^{2} \mathbf{\sigma}^{2} }}{\partial\mathbf{\psi'}\partial\mathbf{\varepsilon}} = \mathbf{B}(m) \, h_{\varepsilon}\left( \mathbf{\varepsilon},z \right)
\label{eq:40}
\end{equation}
\begin{equation}
\frac{{\partial^{2} \mathbf{\sigma}^{2} }}{\partial z\,\partial z'} = \mathbf{GA}_{\sigma}^{-1}\mathbf{AR}_{\sigma}\, h_{zz}\left( \mathbf{\varepsilon},z \right)
\label{eq:41}
\end{equation}
\begin{equation}
\frac{{\partial^{2} \mathbf{\sigma}^{2} }}{\partial z\,\partial\mathbf{\varepsilon}'} = \mathbf{GA}_{\sigma}^{-1}\mathbf{AR}_{\sigma}\, h_{z\varepsilon}\left( \mathbf{\varepsilon},z \right)
\label{eq:42}
\end{equation}
\begin{equation}
\frac{{\partial^{2} \mathbf{\sigma}^{2} }}{\partial\mathbf{\varepsilon}\,\partial\mathbf{\varepsilon}'} = \mathbf{GA}_{\sigma}^{-1}\mathbf{AR}_{\sigma}\, h_{2\varepsilon}\left( \mathbf{\varepsilon},z \right)
\label{eq:H_eps}
\end{equation}
All other second derivatives among $\mathbf{\psi},\mathbf{\omega},\sigma_{0}^{2}$ are zero (since $\mathbf{G}+\mathbf{AR}_{\sigma}h$ is linear in each of these).

For $\boldsymbol{\gamma}$ paired with any of $\mathbf{\psi},\mathbf{\omega},\sigma_{0}^{2},z,\mathbf{\varepsilon}$, only $\mathbf{GA}_{\sigma}^{-1}$ carries $\boldsymbol{\gamma}$-dependence, and the generic notation is genuinely simpler here: writing $\zeta$ for the second parameter,
\begin{equation}
\frac{\partial^{2}\mathbf{\sigma}^{2}}{\partial\boldsymbol{\gamma}\,\partial\zeta'} = \mathbf{B}(k)\cdot\frac{\partial\mathbf{\sigma}^{2}}{\partial\zeta'},
\label{eq:46}
\end{equation}
and for $\boldsymbol{\gamma}$ against itself, differentiating both factors of Equation~\ref{eq:25},
\begin{equation}
\frac{\partial^{2}\mathbf{\sigma}^{2}}{\partial\boldsymbol{\gamma}\,\partial\boldsymbol{\gamma}'} = 2\,\mathbf{B}(k)'\cdot\frac{\partial\mathbf{\sigma}^{2}}{\partial\boldsymbol{\gamma}}.
\label{eq:34}
\end{equation}

\subsection{Derivatives with respect to the ARIMA process}
For $\boldsymbol{\xi},\boldsymbol{\zeta}\in\{\boldsymbol{\rho},\boldsymbol{\lambda},\boldsymbol{\beta}\}$, $\mathbf{\sigma}^{2}$ depends on these parameters only through $\boldsymbol{\varepsilon}(\boldsymbol{\xi})$, so the chain rule applied to Equation~\ref{eq:32} gives
\begin{equation}
\frac{\partial \mathbf{\sigma}^{2} }{\partial \boldsymbol{\xi}} = \mathbf{GA}_{\sigma}^{-1}\mathbf{AR}_{\sigma}\, h_{\varepsilon}\left( \mathbf{\varepsilon}, z \right)\cdot \frac{\partial \mathbf{\varepsilon}}{\partial \boldsymbol{\xi}},
\label{eq:43}
\end{equation}
and, differentiating a second time---an ordinary product rule, since $h_{\varepsilon}$ is itself a function of $\boldsymbol{\varepsilon}(\boldsymbol{\zeta})$---this genuinely benefits from the generic $\boldsymbol{\xi},\boldsymbol{\zeta}$ notation, since it unifies all six $\{\rho,\lambda,\beta\}$-pairs into one equation:
\begin{equation}
\frac{\partial^{2} \mathbf{\sigma}^{2} }{\partial \boldsymbol{\xi}\,\partial\boldsymbol{\zeta}'} =
\mathbf{GA}_{\sigma}^{-1}\mathbf{AR}_{\sigma}\left[\, h_{2\varepsilon}\left( \mathbf{\varepsilon}, z \right)\cdot\frac{\partial \mathbf{\varepsilon}}{\partial \boldsymbol{\xi}}\frac{\partial \mathbf{\varepsilon}}{\partial \boldsymbol{\zeta}'} \;+\; h_{\varepsilon}\left( \mathbf{\varepsilon}, z \right)\cdot\frac{\partial^{2} \mathbf{\varepsilon}}{\partial \boldsymbol{\xi}\,\partial \boldsymbol{\zeta}'}\,\right],
\label{eq:44}
\end{equation}
valid for both the diagonal ($\boldsymbol{\xi}=\boldsymbol{\zeta}$) and off-diagonal case, with $\partial^{2}\boldsymbol{\varepsilon}/\partial\boldsymbol{\xi}\,\partial\boldsymbol{\zeta}'$ given by Equation~\ref{eq:re-two-way-second}. With no or fixed effects this equals the effect transform of Equations~\ref{eq:14}--\ref{eq:17}, which is zero except for pairs involving $\boldsymbol{\lambda}$ and the $\boldsymbol{\rho}$--$\boldsymbol{\beta}$ pair; with random effects all pairs are generally nonzero, because the weights $\theta_j$ depend on all three parameter groups.

For an ARIMA parameter $\boldsymbol{\xi}\in\{\boldsymbol{\rho},\boldsymbol{\lambda},\boldsymbol{\beta}\}$ crossed with a GARCH-native parameter, only $\boldsymbol{\varepsilon}(\boldsymbol{\xi})$ carries the $\boldsymbol{\xi}$-dependence, so each cross term follows directly from the $\mathbf{\varepsilon}$-derivatives already given above:
\begin{equation}
\frac{\partial^{2}\mathbf{\sigma}^{2}}
{\partial \boldsymbol{\xi}\,\partial\boldsymbol{\gamma}'}
=
\mathbf{B}(k)\cdot
\frac{\partial\mathbf{\sigma}^{2}}{\partial\boldsymbol{\xi}},
\tag{47a}
\label{eq:47a}
\end{equation}

\begin{equation}
\frac{\partial^{2}\mathbf{\sigma}^{2}}
{\partial \boldsymbol{\xi}\,\partial\mathbf{\psi}'}
=
\mathbf{B}(m)\,
h_{\varepsilon}\left(\mathbf{\varepsilon},z\right)\cdot
\frac{\partial \mathbf{\varepsilon}}{\partial \boldsymbol{\xi}},
\tag{47b}
\label{eq:47b}
\end{equation}

\begin{equation}
\frac{\partial^{2}\mathbf{\sigma}^{2}}
{\partial \boldsymbol{\xi}\,\partial z'}
=
\mathbf{GA}_{\sigma}^{-1}\mathbf{AR}_{\sigma}\,
h_{z\varepsilon}\left(\mathbf{\varepsilon},z\right)\cdot
\frac{\partial \mathbf{\varepsilon}}{\partial \boldsymbol{\xi}},
\tag{47c}
\label{eq:47c}
\end{equation}
(note, in particular, no leading factor of 2 in the $\boldsymbol{\psi}$ case), while $\partial^{2}\mathbf{\sigma}^{2}/\partial\boldsymbol{\xi}\,\partial\mathbf{\omega}' = \partial^{2}\mathbf{\sigma}^{2}/\partial\boldsymbol{\xi}\,\partial\sigma_{0}^{2\,\prime} = 0$.

Finally, the log-variance derivatives needed for the likelihood (in the section "The LL gradient" below) are worth keeping in generic form, since the same two-line identity applies uniformly to every parameter pair rather than needing to be listed separately for each:
\begin{equation}
\frac{\partial \ln(\mathbf{\sigma}^{2}) }{\partial\xi} = \frac{1}{\mathbf{\sigma}^{2}}\frac{\partial\mathbf{\sigma}^{2}}{\partial\xi},
\label{eq:49}
\end{equation}
\begin{equation}
\frac{\partial^{2}\ln(\boldsymbol{\sigma}^{2})}{\partial\xi\,\partial\zeta'} = \frac{1}{\boldsymbol{\sigma}^{2}}\frac{\partial^{2}\boldsymbol{\sigma}^{2}}{\partial\xi\,\partial\zeta'} - \frac{\partial\ln(\boldsymbol{\sigma}^{2})}{\partial\xi}\frac{\partial\ln(\boldsymbol{\sigma}^{2})}{\partial\zeta'}.
\label{eq:49b}
\end{equation}





\section{Implemented group and time effects}\label{sec:effects-implementation}

The likelihood applies the effects transform to the ARMA innovations $\tilde{\varepsilon}$ of Equation~\ref{eq:arima_reduced}. For an included observation, define the grand mean $\bar{\tilde{\varepsilon}}$, the group mean $\bar{\tilde{\varepsilon}}_{g,i}$, and the mean for its actual calendar date $\bar{\tilde{\varepsilon}}_{t}$. The correction in Equation~\ref{eq:rRE_def} uses a separate weight for each enabled dimension. A `none' effect has weight zero, a `fixed' effect has weight one, and a `random' effect uses the estimated variance components below. The time means are formed from the date map, not from a common matrix-column position.

For the two-way random-effects calculation, let $n_{g,i}=T_i$ be the number of observations in group $i$, and let $n_{t}$ be the number of groups observed at date $t$. Define
\begin{equation}
a_g=\frac{1}{\operatorname{mean}_i(n_{g,i})},\qquad
a_t=\frac{1}{\operatorname{mean}_t(n_t)},\qquad
\mu=\bar{\tilde{\varepsilon}}.
\label{eq:re-count-scales}
\end{equation}
Let $K_s$ be the number of slope regressors, excluding the intercept, and $NT$ the number of included observations. The implementation estimates the idiosyncratic variance from the two-way demeaned innovation:
\begin{equation}
\tilde{\varepsilon}_{w,i,t}=\tilde{\varepsilon}_{i,t}-\bar{\tilde{\varepsilon}}_{g,i}-\bar{\tilde{\varepsilon}}_{t}+\mu,
\qquad
\widehat{\sigma}_{\epsilon}^{2}
=\frac{\operatorname{mean}_{\mathcal O}(\tilde{\varepsilon}_{w,i,t}^{2})}
{1-a_g-a_t+(1-K_s)/NT}.
\label{eq:re-two-way-evar}
\end{equation}
For dimension $j\in\{g,t\}$, first remove the other dimension's mean and add the grand mean; let $\bar{\tilde{\varepsilon}}_{j,\mathrm{adj}}$ denote the resulting mean within the own dimension. Set $n_j^{\mathrm{eff}}=NTa_j$ and $d_j=\max(n_j^{\mathrm{eff}}-1-K_s,1)$. The own-dimension variance estimate is
\begin{equation}
\widehat{\sigma}_{v_j}^{2}
=\max\!\left\{0,
\frac{n_j^{\mathrm{eff}}}{d_j}
\operatorname{mean}_{\mathcal O}\!\left[
(\bar{\tilde{\varepsilon}}_{j,\mathrm{adj}}-\mu)^2\right]
-a_j\widehat{\sigma}_{\epsilon}^{2}\right\}.
\label{eq:re-two-way-vvar}
\end{equation}
The quasi-demeaning weight at an observation is then
\begin{equation}
\theta_{j}=1-\sqrt{\frac{\widehat{\sigma}_{\epsilon}^{2}}
{\widehat{\sigma}_{\epsilon}^{2}+n_{j}\widehat{\sigma}_{v_j}^{2}}},
\label{eq:re-two-way-theta}
\end{equation}
where $n_j$ is that group's observation count for $j=g$, or that date's observed-group count for $j=t$. Negative estimated effect variances are truncated to zero, giving $\theta_j=0$. The variance components are recomputed from the current ARMA innovations during likelihood evaluation; they are not separate entries in the optimizer parameter vector. Since $\tilde{\varepsilon}$ depends on $\boldsymbol{\beta}$, $\boldsymbol{\rho}$ and $\boldsymbol{\lambda}$, so do the weights. The likelihood applies group and time effects to the innovations of the mean equation; the `effects.variance' setting does not add a variance-effects transform.

\subsection{First and second derivatives of the implemented two-way effects}
The following derivatives describe the effect transform of the ARMA innovations $\tilde{\varepsilon}$. Let $\langle x\rangle=(NT)^{-1}\sum_{\mathcal O}x$ be the mean over included observations, let $\bar{x}_{j}$ be the mean within the observation's own group ($j=g$) or actual calendar date ($j=t$), and define the centered operator
\begin{equation}
\mathcal C_j(x)=\langle x\rangle-\bar{x}_{j},\qquad
\varepsilon=\tilde{\varepsilon}+\sum_{j\in\{g,t\}}\theta_j\mathcal C_j(\tilde{\varepsilon}).
\label{eq:re-two-way-transform}
\end{equation}
For parameters $h,g\in\{\boldsymbol{\beta},\boldsymbol{\rho},\boldsymbol{\lambda}\}$, write $\tilde{\varepsilon}_h=\partial\tilde{\varepsilon}/\partial h$ and $\tilde{\varepsilon}_{hg}=\partial^2\tilde{\varepsilon}/\partial h\,\partial g$, given by Equations~\ref{eq:11}--\ref{eq:17}. For any such derivative $x\in\{\tilde{\varepsilon}_h,\tilde{\varepsilon}_{hg}\}$, define the linear maps
\begin{align}
w(x)&=x-\bar{x}_{g}-\bar{x}_{t}+\langle x\rangle,\\
b_{j}(x)&=\overline{\left(x-\bar{x}_{-j}+\langle x\rangle\right)}_{j}-\langle x\rangle,
\label{eq:re-two-way-input-derivatives}
\end{align}
where $\bar{x}_{-j}$ is the mean over the other panel dimension, and the overbar outside parentheses denotes the mean over dimension $j$. Write $\tilde{\varepsilon}_{w,h}=w(\tilde{\varepsilon}_h)$, $b_{j,h}=b_j(\tilde{\varepsilon}_h)$, $\tilde{\varepsilon}_{w,hg}=w(\tilde{\varepsilon}_{hg})$ and $b_{j,hg}=b_j(\tilde{\varepsilon}_{hg})$. Let $D=1-a_g-a_t+(1-K_s)/NT$ be the denominator in Equation~\ref{eq:re-two-way-evar}, $n_j^{\mathrm{eff}}=NTa_j$, $d_j=\max(n_j^{\mathrm{eff}}-1-K_s,1)$, and $c_j=n_j^{\mathrm{eff}}/(d_jNT)$. On the region where the variance components are positive, their first derivatives are
\begin{align}
\tau_h
&\equiv\frac{\partial\widehat{\sigma}_{\epsilon}^{2}}{\partial h}
=\frac{2}{NTD}\sum_{\mathcal O}\tilde{\varepsilon}_w\tilde{\varepsilon}_{w,h},
\label{eq:re-two-way-tau-first}\\
\nu_{j,h}
&\equiv\frac{\partial\widehat{\sigma}_{v_j}^{2}}{\partial h}
=2c_j\sum_{\mathcal O}(\bar{\tilde{\varepsilon}}_{j,\mathrm{adj}}-\mu)b_{j,h}-a_j\tau_h.
\label{eq:re-two-way-nu-first}
\end{align}
Here $\tau=\widehat{\sigma}_{\epsilon}^{2}$ and $\nu_j=\widehat{\sigma}_{v_j}^{2}$. For an observation with own-dimension count $n_j$, the weight derivatives are
\begin{align}
\theta_{j,\tau}&=-\frac{1}{2\tau}(1-\theta_j)\theta_j(2-\theta_j),\\
\theta_{j,\nu}&=\frac{n_j}{2\tau}(1-\theta_j)^3,\\
\theta_{j,h}&=\theta_{j,\tau}\tau_h+\theta_{j,\nu}\nu_{j,h}.
\label{eq:re-two-way-theta-first}
\end{align}
Thus the first derivative of the transformed innovation is
\begin{equation}
\varepsilon_{h}=\tilde{\varepsilon}_h+\sum_{j\in\{g,t\}}\left[\theta_j\mathcal C_j(\tilde{\varepsilon}_h)+\theta_{j,h}\mathcal C_j(\tilde{\varepsilon})\right].
\label{eq:re-two-way-first}
\end{equation}

Unlike the regression residual $u$, the innovation $\tilde{\varepsilon}$ is not linear in the parameters, so $\tilde{\varepsilon}_{hg}$ is generally nonzero (Equations~\ref{eq:14}--\ref{eq:17}). The second derivatives of the variance components are therefore
\begin{equation}
\tau_{hg}=\frac{2}{NTD}\sum_{\mathcal O}\left(\tilde{\varepsilon}_{w,h}\tilde{\varepsilon}_{w,g}+\tilde{\varepsilon}_{w}\tilde{\varepsilon}_{w,hg}\right),\qquad
\nu_{j,hg}=2c_j\sum_{\mathcal O}\left[b_{j,h}b_{j,g}+(\bar{\tilde{\varepsilon}}_{j,\mathrm{adj}}-\mu)b_{j,hg}\right]-a_j\tau_{hg}.
\label{eq:re-two-way-variance-second}
\end{equation}
The second derivatives of each weight follow from $\theta_j=\theta_j(\tau,\nu_j)$:
\begin{align}
\theta_{j,\tau\tau}
&=-\frac{\theta_{j,\tau}}{2\tau}\left[4-3(2-\theta_j)\theta_j\right],\\
\theta_{j,\nu\nu}
&=-\frac{3}{4}\left(\frac{n_j}{\tau}\right)^2(1-\theta_j)^5,\\
\theta_{j,\tau\nu}
&=-\frac{\theta_{j,\nu}}{2\tau}\left[3(\theta_j-2)\theta_j+2\right],\\
\theta_{j,hg}
&=\theta_{j,\tau\tau}\tau_h\tau_g
+\theta_{j,\tau\nu}(\tau_h\nu_{j,g}+\nu_{j,h}\tau_g)
+\theta_{j,\nu\nu}\nu_{j,h}\nu_{j,g}
+\theta_{j,\tau}\tau_{hg}+\theta_{j,\nu}\nu_{j,hg}.
\label{eq:re-two-way-theta-second}
\end{align}
The corresponding second derivative of the transformed innovation is
\begin{equation}
\varepsilon_{hg}=\tilde{\varepsilon}_{hg}+\sum_{j\in\{g,t\}}\left[\theta_j\mathcal C_j(\tilde{\varepsilon}_{hg})+\theta_{j,h}\mathcal C_j(\tilde{\varepsilon}_g)+\theta_{j,g}\mathcal C_j(\tilde{\varepsilon}_h)+\theta_{j,hg}\mathcal C_j(\tilde{\varepsilon})\right].
\label{eq:re-two-way-second}
\end{equation}
These derivatives apply to a random-effects dimension while $\tau>0$, $\nu_j>0$, and $n_j>1$. For a fixed effect, $\theta_j=1$ is constant; for a disabled effect, $\theta_j=0$, so its weight derivatives vanish and Equations~\ref{eq:re-two-way-first} and~\ref{eq:re-two-way-second} reduce to the linear transform of $\tilde{\varepsilon}_h$ and $\tilde{\varepsilon}_{hg}$. The implementation also sets the random-effects weight and its derivatives to zero when the estimated own-dimension variance is clipped at zero or the group/date has at most one observation.

\subsection{One-way random-effects reference (not the implementation)}

For comparison, a textbook one-way random-effects model can be written as

\begin{equation}
\eta_{i,t} = v_{i} + \epsilon_{i,t}
\label{eq:50}
\end{equation}

The equations that follow in this subsection give a one-way Swamy--Arora-style reference model. They are not the variance-component calculation used by the current package, which is the two-way procedure in Equations~\ref{eq:re-count-scales}--\ref{eq:re-two-way-theta} and works on regression residuals before ARMA/GARCH standardization.

In that one-way reference model, normalized residuals could be used as:

\begin{equation}
\eta_{i,t} = \frac{\varepsilon_{i,t}}{\sigma_{i,t}}
\label{eq:51}
\end{equation}

With known parameters, but unknown group effects, the idiosyncratic
residuals can be obtained using the Swamy--Arora-type variance-components weight \citep{swamy1972exact}:

\begin{equation}
\theta_{i} = 1 - \frac{\sigma_{\epsilon}}{\sqrt{\sigma_{\epsilon}^{2} + T_{i}\sigma_{v}^{2}}}
\label{eq:re1}
\end{equation}

And then \citep[pp.~412--413]{greene2012econometric}, using the
group means ${\overline{\eta}}_{i}$ to calculate

\begin{equation}
\epsilon_{i,t} = \eta_{i,t} + RE_{0}(\eta_{i,t},\theta_{i})\
\label{eq:re2}
\end{equation}

Where the random effects function, indexed by $j \in \{0,1\}$, is

\begin{equation}
RE_{j}\left( x_{i,t},\theta_{i} \right) = j \cdot x_{i,t} + \theta_{i}\left( \overline{x} - {\overline{x}}_{i} \right)
\label{eq:re3}
\end{equation}
(note that the subscript $j$ here denotes the weight on $x_{i,t}$ itself, and should not be confused with the group index $i$),
and ${\overline{x}}_{i}$ are the group means of $x_{i,t}$ and
$\overline{x}$ the mean of ${\overline{x}}_{i}$, ensuring a centered
effect. $\sigma_{\epsilon}^{2}$ is estimated as

\begin{equation}{\widehat{\sigma}}_{\epsilon}^{2} = \frac{1}{1 - \overline{1/T}}\frac{1}{N\sum_{}^{}T_{i}}\sum_{i = 0}^{N}{\sum_{t = 0}^{T_{i}}\left( RE_{1}\left( \eta_{i,t},1 \right) \right)^{2}}
\label{eq:re4}
\end{equation}

Where

\begin{equation}
\overline{1/T} = \frac{1}{N}\sum_{i = 0}^{N}\frac{1}{T}
\label{eq:re5}
\end{equation}

And $\sigma_{v}^{2}$ is estimated as:

\begin{equation}
{\widehat{\sigma}}_{v}^{2} = \frac{1}{N\sum_{}^{}T_{i}}\sum_{i = 1}^{N}{\sum_{t = 0}^{T_{i}}\eta_{i,t}^{2}}
\label{eq:re6}
\end{equation}

The $\theta_{i}$ derivatives are

\begin{equation}
\frac{\partial\theta_{i}}{\partial\sigma_{v}^{2}} = \frac{1}{2}\frac{T_{i}}{\sigma_{\epsilon}^{2}}\left( 1 - \theta_{i} \right)^{3}
\label{eq:re7}
\end{equation}

\begin{equation}
\frac{\partial\theta_{i}}{\partial\sigma_{\epsilon}^{2}} = - \frac{1}{2}\frac{1}{\sigma_{\epsilon}^{2}}\left( 1 - \theta_{i} \right)\theta_{i}\left( 2 - \theta_{i} \right)
\label{eq:re7b}
\end{equation}

\begin{equation}
\frac{\partial^{2}\theta_{i}}{{\partial\left( \sigma_{\epsilon}^{2} \right)}^{2}} = - \frac{\partial\theta_{i}}{\partial\sigma_{\epsilon}^{2}}\frac{1}{2}\frac{4 - 3\left( 2 - \theta_{i} \right)\theta_{i}}{\sigma_{\epsilon}^{2}}
\label{eq:re8}
\end{equation}

\begin{equation}
\frac{\partial\theta_{i}}{\partial\left( \sigma_{v}^{2} \right)^{2}} = - \frac{3}{4}\frac{T_{i}^{2}}{\left( \sigma_{\epsilon}^{2} \right)^{2}}\left( 1 - \theta_{i} \right)^{5}
\label{eq:re9}
\end{equation}

\begin{equation}\frac{\partial\theta_{i}}{\partial\sigma_{v}^{2}\partial\sigma_{\epsilon}^{2}} = - \frac{\partial\epsilon_{i,t}}{\partial\sigma_{v}^{2}}\frac{1}{2}\frac{3\left( \theta_{i} - 2 \right)\theta_{i} - 2}{\sigma_{\epsilon}^{2}}
\label{eq:re10}
\end{equation}

Now, for any parameter $h$ that affects the original residuals, the
marginal effect on the random effect residuals (recall Equation~\ref{eq:re2}) is

\begin{align}
\frac{\partial \epsilon_{i,t}}{\partial h} 
= {} & \frac{\partial \eta_{i,t}}{\partial h} 
+ \widehat{\theta}_{i} \left( \frac{\partial \overline{\eta}}{\partial h} 
- \frac{\partial \overline{\eta}_{i}}{\partial h} \right) 
+ \frac{\partial \widehat{\theta}_{i}}{\partial h} 
\left( \overline{\eta} - \overline{\eta}_{i} \right) \nonumber \\
= {} & \frac{\partial \eta_{i,t}}{\partial h} 
+ RE_{0} \left( \frac{\partial \eta_{i,t}}{\partial h}, \widehat{\theta}_{i} \right) 
+ RE_{0} \left( \eta_{i,t}, \frac{\partial \widehat{\theta}_{i}}{\partial h} \right)
\label{eq:re12}
\end{align}


Where

\begin{equation}
\frac{\partial{\widehat{\theta}}_{i}}{\partial h} = - \frac{1}{2}\sqrt{\frac{1}{{\widehat{\sigma}}_{\epsilon}^{2}{\widehat{\sigma}}_{z}^{2}\left( T_{i} - 1 \right)}}\frac{\partial{\widehat{\sigma}}_{\epsilon}^{2}}{\partial h}
\label{eq:re13}
\end{equation}

\begin{align}
\frac{\partial \widehat{\sigma}_{\epsilon}^{2}}{\partial h} 
= {} & 2 \sum_{i = 1}^{N} \sum_{t = 1}^{T} 
RE_{1}\left( \eta_{i,t}, 1 \right) \left( 
\frac{\partial \eta_{i,t}}{\partial h} 
+ \frac{1}{NT} \sum_{i=1}^{N} \sum_{t = 1}^{T} \frac{\partial \eta_{i,t}}{\partial h} 
- \frac{1}{T_{i}} \sum_{t = 1}^{T_{i}} \frac{\partial \eta_{i,t}}{\partial h} 
\right) \nonumber \\
= {} & \sum_{i = 1}^{N} \sum_{t = 1}^{T} 
RE_{1}\left( \eta_{i,t}, 1 \right) 
RE_{1}\left( \frac{\partial \eta_{i,t}}{\partial h}, 1 \right)
\label{eq:re14}
\end{align}

And the second derivative with respect to $h$ and $g$ is

\begin{equation}
\frac{\partial^{2}\epsilon_{i,t}}{\partial g\partial h} = {\widehat{\theta}}_{i}\left( \frac{\partial^{2}\overline{\varepsilon}}{\partial h\partial g} - \frac{\partial^{2}{\overline{\varepsilon}}_{i}}{\partial h\partial g} \right) - \frac{\partial{\widehat{\theta}}_{i}}{\partial h}\frac{\partial{\overline{\varepsilon}}_{i}}{\partial g} - \frac{\partial{\widehat{\theta}}_{i}}{\partial g}\frac{\partial{\overline{\varepsilon}}_{i}}{\partial h} - \frac{\partial^{2}{\widehat{\theta}}_{i}}{\partial h\partial g}{\overline{\varepsilon}}_{i}
\label{eq:re15}
\end{equation}

\begin{align}
\frac{\partial^{2} \epsilon_{i,t}}{\partial g \, \partial h} = {} & 
\frac{\partial^{2} \eta_{i,t}}{\partial g \, \partial h} 
+ RE_{0} \left( \frac{\partial^{2} \eta_{i,t}}{\partial g \, \partial h}, \widehat{\theta}_{i} \right) 
+ RE_{0} \left( \frac{\partial \eta_{i,t}}{\partial g}, \frac{\partial \widehat{\theta}_{i}}{\partial h} \right) \nonumber \\
& + RE_{0} \left( \frac{\partial \eta_{i,t}}{\partial h}, \frac{\partial \widehat{\theta}_{i}}{\partial g} \right) 
+ RE_{0} \left( \eta_{i,t}, \frac{\partial^{2} \widehat{\theta}_{i}}{\partial h \, \partial g} \right)
\label{eq:re16}
\end{align}


Where

\begin{equation}
\frac{\partial^{2}{\widehat{\theta}}_{i}}{\partial h\partial g} = \sqrt{\frac{1}{{\widehat{\sigma}}_{\epsilon}^{2}{\widehat{\sigma}}_{z}^{2}\left( T_{i} - 1 \right)}}\left( \frac{1}{4}\frac{1}{{\widehat{\sigma}}_{\epsilon}^{2}}\frac{\partial{\widehat{\sigma}}_{\epsilon}^{2}}{\partial h}\frac{\partial{\widehat{\sigma}}_{\epsilon}^{2}}{\partial g} - \frac{1}{2}\frac{\partial^{2}{\widehat{\sigma}}_{\epsilon}^{2}}{\partial h\partial g} \right)
\label{eq:re17}
\end{equation}

\begin{align}
\frac{\partial^{2} \widehat{\sigma}_{\epsilon}^{2}}{\partial h \, \partial g} = {} & 
2 \sum_{i = 1}^{N} \sum_{t = 1}^{T} 
RE_{1} \left( \frac{\partial \eta_{i,t}}{\partial h}, 1 \right) 
RE_{1} \left( \frac{\partial \eta_{i,t}}{\partial g}, 1 \right) \nonumber \\
& + 2 \sum_{i = 1}^{N} \sum_{t = 1}^{T} 
RE_{1} \left( \eta_{i,t}, 1 \right) 
RE_{1} \left( \frac{\partial^{2} \widehat{\theta}_{i}}{\partial h \, \partial g}, 1 \right)
\label{eq:re18}
\end{align}

\section{The LL gradient}\label{the-ll-gradient}

Given the normal distribution in Equation~\ref{eq:normal_pdf}, it follows that the log-likelihood contribution of group $i$, summed over its $T_i$ time periods, is

\begin{equation}
{LL}_{i} = - \frac{1}{2}\ln(2\pi) \cdot T_i - \frac{1}{2}\mathbf{1}^{'}\left(\ln\left(\mathbf{\sigma}_{i,t}^{2} \right)+\frac{{\mathbf{\epsilon}_{i,t}^{2}}}{\mathbf{\sigma}_{i,t}^{2}}\right)
\label{eq:52}
\end{equation}

so that, pointwise in $t$,

\begin{equation}
\frac{\partial LL}{\partial\epsilon_{i,t}} = - \frac{{\mathbf{\epsilon}_{i,t}}}{\mathbf{\sigma}_{i,t}^{2}}
\label{eq:55}
\end{equation}

and

\begin{equation}
\frac{\partial LL}{\partial\ln\left(\mathbf{\sigma}_{i,t}^{2}\right)} = - \frac{1}{2}\left(1-\frac{{\mathbf{\epsilon}_{i,t}^{2}}}{\mathbf{\sigma}_{i,t}^{2}}\right)
\label{eq:56}
\end{equation}

The likelihood code exposes derivatives with respect to the variance state used by each model. For ordinary GARCH, this state is $v=\sigma^{2}$; for EGARCH, it is the log-variance $r=\ln v$. The log-variance derivatives above convert to variance derivatives as
\begin{equation}
\frac{\partial LL}{\partial v}=\frac{1}{v}\frac{\partial LL}{\partial r},\qquad
\frac{\partial^{2}LL}{\partial\epsilon\,\partial v}=\frac{1}{v}\frac{\partial^{2}LL}{\partial\epsilon\,\partial r},\qquad
\frac{\partial^{2}LL}{\partial v^{2}}=\frac{1}{v^{2}}\left(\frac{\partial^{2}LL}{\partial r^{2}}-\frac{\partial LL}{\partial r}\right).
\label{eq:variance-logvariance-chain}
\end{equation}

\subsection{Adjusted Gaussian likelihood}
When \texttt{GARCH\_assist} or \texttt{kurtosis\_adj} is nonzero, the implementation selects its \texttt{Hyperbolic} likelihood instead of the default Gaussian likelihood. Let $q=\epsilon^{2}+10^{-8}$, $s=\operatorname{sign}(q-v)$, $a=\texttt{GARCH\_assist}$, $k=\texttt{kurtosis\_adj}$, and let $m_v$ denote the variance-boundary mask (zero at an active bound and one in the interior). The implemented pointwise likelihood is
\begin{equation}
L_H=-\frac{1}{2}\ln(2\pi)-\frac{1}{2}\left[\ln v+(1-k)\frac{q}{v}+a\frac{|q-v|}{v}+\frac{k}{3}\frac{q^{2}}{v^{2}}\right].
\label{eq:hyperbolic-ll}
\end{equation}
Away from the kink $q=v$, its score and Hessian entries are
\begin{align}
\frac{\partial L_H}{\partial\epsilon}
&=-\frac{1}{2}\left[2(1-k)\frac{\epsilon}{v}+2as\frac{\epsilon}{v}+\frac{4k}{3}\frac{q\epsilon}{v^{2}}\right],
\label{eq:hyperbolic-score-e}\\
\frac{\partial L_H}{\partial v}
&=-\frac{1}{2}m_v\left[\frac{1}{v}-(1-k)\frac{q}{v^{2}}-a\frac{s}{v}-a\frac{|q-v|}{v^{2}}-\frac{2k}{3}\frac{q^{2}}{v^{3}}\right],
\label{eq:hyperbolic-score-v}\\
\frac{\partial^{2}L_H}{\partial\epsilon^{2}}
&=-\frac{1}{2}\left[\frac{2(1-k)}{v}+\frac{2as}{v}+\frac{4k}{3}\frac{q+2\epsilon^{2}}{v^{2}}\right],
\label{eq:hyperbolic-hess-ee}\\
\frac{\partial^{2}L_H}{\partial v\,\partial\epsilon}
&=-\frac{1}{2}m_v\left[-2(1-k)\frac{\epsilon}{v^{2}}-2as\frac{\epsilon}{v^{2}}-\frac{8k}{3}\frac{q\epsilon}{v^{3}}\right],
\label{eq:hyperbolic-hess-ve}\\
\frac{\partial^{2}L_H}{\partial v^{2}}
&=-\frac{1}{2}m_v\left[-\frac{1}{v^{2}}+2(1-k)\frac{q}{v^{3}}+2a\frac{s}{v^{2}}+2a\frac{|q-v|}{v^{3}}+2k\frac{q^{2}}{v^{4}}\right].
\label{eq:hyperbolic-hess-vv}
\end{align}
At $q=v$, the absolute-value term is not differentiable; the implementation uses $\operatorname{sign}(0)=0$ there.


\begin{equation}
\frac{\partial LL}{\partial\mathbf{\xi}} = \mathbf{1}^{'}\left(\frac{dLL}{d\epsilon_{i,t}}\frac{\partial\epsilon_{i,t}}{\partial\mathbf{\xi}} + \frac{\partial LL}{{\partial ln}\left( \sigma_{i,t}^{2} \right)}\frac{{\partial ln}\left( \mathbf{\sigma}_{i,t}^{2} \right)}{\partial\mathbf{\xi}} \right)
\label{eq:57}
\end{equation}

For
$\mathbf{\xi} \in \left\{ \mathbf{\beta}\mathbf{,}\mathbf{\rho}\mathbf{,}\mathbf{\lambda}\mathbf{,}\mathbf{\gamma}\mathbf{,}\mathbf{\psi}\mathbf{,}\mathbf{\omega}\mathbf{,}z \right\}$.
The first term is zero for
$\mathbf{\xi} \in \left\{ \mathbf{\psi,\omega,}z \right\}$ since
these parameters do not affect the residuals.



\section{The Hessian}\label{the-hessian}



\begin{equation}
\frac{\partial^{2}dLL}{\partial\mathbf{\varepsilon}_{i}^{2}}=
- \mathbf{1}'\frac{1}{\mathbf{\sigma}_{i}^{2}}
\label{eq:H1}
\end{equation}


\begin{equation}
\frac{\partial^{2}dLL}{{\partial ln}\left( \mathbf{\sigma}_{i}^{2} \right)\partial\mathbf{\varepsilon}_{i}} = \mathbf{1}'\frac{\mathbf{\epsilon}_{i}}{\mathbf{\sigma}_{i}^{2}}
\label{eq:H2}
\end{equation}


\begin{equation}\frac{\partial^{2}dLL}{{\partial ln}\left( \mathbf{\sigma}_{i}^{2} \right)^{2}} = - \frac{1}{2}\frac{{\mathbf{\epsilon}_{i}}^{2}}{\mathbf{\sigma}_{i}^{2}}
\label{eq:H3}
\end{equation}




For parameter $\mathbf{\xi}$ and $\mathbf{\zeta}$, the corresponding hessian elements are


\begin{align}
\frac{\partial^{2} dLL}{\partial \boldsymbol{\xi} \, \partial \boldsymbol{\zeta}'} = {} & 
\mathbf{1}' \Bigg(
\frac{\partial^{2} LL}{\partial \boldsymbol{\epsilon}_{i}^{2}} 
\cdot \frac{\partial \boldsymbol{\epsilon}_{i}}{\partial \boldsymbol{\xi}} 
\cdot \frac{\partial \boldsymbol{\epsilon}_{i}}{\partial \boldsymbol{\zeta}'} \nonumber \\
& + \frac{\partial^{2} dLL}{\partial \ln\left( \boldsymbol{\sigma}_{i}^{2} \right) \, \partial \boldsymbol{\epsilon}_{i}} 
\cdot \left(
\frac{\partial \boldsymbol{\epsilon}_{i}}{\partial \boldsymbol{\xi}} 
\cdot \frac{\partial \ln\left( \boldsymbol{\sigma}_{i}^{2} \right)}{\partial \boldsymbol{\zeta}'} 
+ \frac{\partial \ln\left( \boldsymbol{\sigma}_{i}^{2} \right)}{\partial \boldsymbol{\xi}} 
\cdot \frac{\partial \boldsymbol{\epsilon}_{i}}{\partial \boldsymbol{\zeta}'} 
\right) \nonumber \\
& + \frac{\partial^{2} dLL}{\partial \ln\left( \boldsymbol{\sigma}_{i}^{2} \right)^{2}} 
\cdot \frac{\partial \ln\left( \boldsymbol{\sigma}_{i}^{2} \right)}{\partial \boldsymbol{\xi}} 
\cdot \frac{\partial \ln\left( \boldsymbol{\sigma}_{i}^{2} \right)}{\partial \boldsymbol{\zeta}'} \nonumber \\
& + \frac{\partial LL}{\partial \boldsymbol{\epsilon}_{i}} 
\cdot \frac{\partial^{2} \boldsymbol{\epsilon}_{i}}{\partial \boldsymbol{\xi} \, \partial \boldsymbol{\zeta}'} 
+ \frac{\partial LL}{\partial \ln\left( \boldsymbol{\sigma}_{i}^{2} \right)} 
\cdot \frac{\partial^{2} \ln\left( \boldsymbol{\sigma}_{i}^{2} \right)}{\partial \boldsymbol{\xi} \, \partial \boldsymbol{\zeta}'} 
\Bigg)
\label{eq:H4}
\end{align}


Where $\frac{\partial\mathbf{\varepsilon}_{i}}{\partial\mathbf{\xi}} = \frac{\partial^{2}\mathbf{\varepsilon}_{i}}{\partial\mathbf{\xi}\partial\mathbf{\zeta'}} = 0$ for $\mathbf{\gamma,\ \psi}$ and $\mathbf{\omega}$. The hessian is
then


\begin{tabular}{llllllllll}
$\frac{\partial^{2} LL}{\partial \boldsymbol{\beta} \, \partial \boldsymbol{\beta}'}$ & \\

$\frac{\partial^{2} LL}{\partial \boldsymbol{\rho} \, \partial \boldsymbol{\beta}'}$ &
$\frac{\partial^{2} LL}{\partial \boldsymbol{\rho} \, \partial \boldsymbol{\rho}'}$ & & & & & & & \\

$\frac{\partial^{2} LL}{\partial \boldsymbol{\lambda} \, \partial \boldsymbol{\beta}'}$ &
$\frac{\partial^{2} LL}{\partial \boldsymbol{\lambda} \, \partial \boldsymbol{\rho}'}$ &
$\frac{\partial^{2} LL}{\partial \boldsymbol{\lambda} \, \partial \boldsymbol{\lambda}'}$ & & & & & & & \\

$\frac{\partial^{2} LL}{\partial \boldsymbol{\gamma} \, \partial \boldsymbol{\beta}'}$ &
$\frac{\partial^{2} LL}{\partial \boldsymbol{\gamma} \, \partial \boldsymbol{\rho}'}$ &
$\frac{\partial^{2} LL}{\partial \boldsymbol{\gamma} \, \partial \boldsymbol{\lambda}'}$ &
$\frac{\partial^{2} LL}{\partial \boldsymbol{\gamma} \, \partial \boldsymbol{\gamma}'}$ & & & & & & \\

$\frac{\partial^{2} LL}{\partial \boldsymbol{\psi} \, \partial \boldsymbol{\beta}'}$ &
$\frac{\partial^{2} LL}{\partial \boldsymbol{\psi} \, \partial \boldsymbol{\rho}'}$ &
$\frac{\partial^{2} LL}{\partial \boldsymbol{\psi} \, \partial \boldsymbol{\lambda}'}$ &
$\frac{\partial^{2} LL}{\partial \boldsymbol{\psi} \, \partial \boldsymbol{\gamma}'}$ &
$\frac{\partial^{2} LL}{\partial \boldsymbol{\psi} \, \partial \boldsymbol{\psi}'}$ & & & & & \\

$\frac{\partial^{2} LL}{\partial \boldsymbol{\omega} \, \partial \boldsymbol{\beta}'}$ &
$\frac{\partial^{2} LL}{\partial \boldsymbol{\omega} \, \partial \boldsymbol{\rho}'}$ &
$\frac{\partial^{2} LL}{\partial \boldsymbol{\omega} \, \partial \boldsymbol{\lambda}'}$ &
$\frac{\partial^{2} LL}{\partial \boldsymbol{\omega} \, \partial \boldsymbol{\gamma}'}$ &
$\frac{\partial^{2} LL}{\partial \boldsymbol{\omega} \, \partial \boldsymbol{\psi}'}$ &
$\frac{\partial^{2} LL}{\partial \boldsymbol{\omega} \, \partial \boldsymbol{\omega}'}$ & & & & \\

$\frac{\partial^{2} LL}{\partial z \, \partial \boldsymbol{\beta}'}$ &
$\frac{\partial^{2} LL}{\partial z \, \partial \boldsymbol{\rho}'}$ &
$\frac{\partial^{2} LL}{\partial z \, \partial \boldsymbol{\lambda}'}$ &
$\frac{\partial^{2} LL}{\partial z \, \partial \boldsymbol{\gamma}'}$ &
$\frac{\partial^{2} LL}{\partial z \, \partial \boldsymbol{\psi}'}$ &
$\frac{\partial^{2} LL}{\partial z \, \partial \boldsymbol{\omega}'}$ &
$\frac{\partial^{2} LL}{\partial z \partial z'}$ & & \\


\end{tabular}

\section{Data-Free Expected Standard Errors}\label{sec:fisher-information}

\paragraph{AI-generation disclosure.}
Unlike my original derivations elsewhere in this document, this section was generated entirely by AI. 
The reason is twofold; first, the asymptotica approximation is for evaluation purposes only, not for estimation. hence this math is not 
part of the package and does not affect results. Second, obtaining data free expectations of standard errors for this all ready 
very complicated model, is very challenging, and given that it is not part of the package, not worth the effort to obtain it rigorously. 

\subsection{Target}
For simulation studies it is useful to have a reference standard error that is computed from the data-generating process alone, without simulating or observing any data. We use the inverse expected information at the true parameters $\boldsymbol\vartheta_0$,
\begin{equation}
\operatorname{SE}_j=\sqrt{\big[\mathcal I^{-1}\big]_{jj}},\qquad
\mathcal I=\mathbb E_{\boldsymbol\vartheta_0}\!\left[-\mathbf H(\boldsymbol\vartheta_0)\right],
\label{eq:fisher-definition}
\end{equation}
for the panel dimensions $N$ and $T$ of the design. This is the curvature reference that the reported (Hessian-based) standard errors estimate. It is not the sampling standard deviation of the estimates. When the likelihood is strongly non-quadratic, as for the GARCH variance constant whose estimate is almost perfectly negatively correlated with the estimated persistence $\sum\psi_m+\sum\gamma_k$, the sampling standard deviation and the standard errors evaluated at the estimates can both exceed Equation~\ref{eq:fisher-definition} in moderate samples.

The calculation evaluates $\mathcal I$ exactly up to one approximation, the replacement of the random conditional variance by its expectation (Section~\ref{sec:ratio-approx}), and two finite-sample simplifications: the random-effect weights are held at their expected values, and the conditional-innovation property is assumed for the effect-transformed residual. The approximation does not vanish as the sample grows. The simplifications are of order $1/N$ and $1/T$.

\subsection{Design}
The data-generating process is, for group $i=1,\dots,N$ and $t=1,\dots,T$,
\begin{equation}
y_{it}=\beta_0+\mathbf{x}_{it}'\boldsymbol{\beta}_1+\alpha_i+\lambda_t+u_{it},\qquad
\rho(\mathbf{L})u_{it}=\lambda(\mathbf{L})\epsilon_{it},\qquad
\epsilon_{it}=\sigma_{it}z_{it},
\end{equation}
with $z_{it}$ i.i.d.\ standard normal, $\sigma_{it}^2=\omega+\sum_m\psi_m\epsilon_{i,t-m}^2+\sum_k\gamma_k\sigma_{i,t-k}^2$ stationary, $\alpha_i\sim\mathcal N(0,\sigma_\alpha^2)$, $\lambda_t\sim\mathcal N(0,\sigma_\lambda^2)$, and each regressor an independent ARMA--GARCH process with the same dynamics and innovation variance $\sigma_x^2$. We assume stationary and invertible ARMA polynomials, $\sum_m\psi_m+\sum_k\gamma_k<1$, a finite fourth moment of $\epsilon_{it}$ (Equation~\ref{eq:fourth-moment-condition}), at least one $\psi_m\neq0$ when GARCH terms are fitted, fitted orders at least as large as the true orders, a balanced panel, and no differencing.

Let $\boldsymbol\Gamma$ be the $T\times T$ Toeplitz autocovariance matrix of $u_{it}$ (from the MA($\infty$) weights of $\lambda(\mathbf L)/\rho(\mathbf L)$ and $\sigma_\epsilon^2=\omega/(1-\sum\psi_m-\sum\gamma_k)$), $\boldsymbol\Gamma_x$ the same matrix for a regressor, $\mathbf 1$ a vector of ones and $\mathbf P=T^{-1}\mathbf 1\mathbf 1'$.

\subsection{The filtered and transformed residual}
At the true parameters the regression residual of group $i$ is $u_i+\alpha_i\mathbf 1+\boldsymbol\lambda$. With $\mathbf M_0=\mathbf{MA}^{-1}\mathbf{AR}$ the lower-triangular ARMA filter of Equation~\ref{eq:arima_reduced} (zero presample), the ARMA innovation is $\tilde{\boldsymbol\varepsilon}_i=\mathbf M_0(\mathbf u_i+\alpha_i\mathbf 1+\boldsymbol\lambda)$. Its covariance splits into a part that is independent across groups and a part common to all groups:
\begin{equation}
\boldsymbol\Omega=\mathbf M_0\big(\boldsymbol\Gamma+\sigma_\alpha^2\mathbf 1\mathbf 1'\big)\mathbf M_0',\qquad
\boldsymbol\Lambda=\sigma_\lambda^2\mathbf M_0\mathbf M_0'.
\label{eq:df-omega-lambda}
\end{equation}
In a balanced panel the two-way transform of Equation~\ref{eq:rRE_def} is linear,
\begin{equation}
\boldsymbol\varepsilon_i=\mathbf R_{\mathrm{own}}\tilde{\boldsymbol\varepsilon}_i+\mathbf R_{\mathrm{oth}}\sum_{j\neq i}\tilde{\boldsymbol\varepsilon}_j,\qquad
\begin{aligned}
\mathbf R_{\mathrm{own}}&=\big(1-\tfrac{\theta_t}{N}\big)\mathbf I-\big(\theta_g-\tfrac{\theta_g+\theta_t}{N}\big)\mathbf P,\\
\mathbf R_{\mathrm{oth}}&=-\tfrac{\theta_t}{N}\mathbf I+\tfrac{\theta_g+\theta_t}{N}\mathbf P,
\end{aligned}
\label{eq:df-transform}
\end{equation}
and a component common to all groups is multiplied by $\mathbf S=\mathbf R_{\mathrm{own}}+(N-1)\mathbf R_{\mathrm{oth}}=(1-\theta_t)\mathbf I+\theta_t\mathbf P$. For any two linear maps $\mathbf F_a,\mathbf F_b$ of the residual define
\begin{equation}
\begin{aligned}
\mathcal C(\mathbf F_a,\mathbf F_b)={}&\mathbf R_{\mathrm{own}}\mathbf F_a\boldsymbol\Sigma\mathbf F_b'\mathbf R_{\mathrm{own}}'+(N-1)\mathbf R_{\mathrm{oth}}\mathbf F_a\boldsymbol\Sigma\mathbf F_b'\mathbf R_{\mathrm{oth}}'\\
&+\sigma_\lambda^2\mathbf S\mathbf F_a\mathbf F_b'\mathbf S',
\qquad \boldsymbol\Sigma=\boldsymbol\Gamma+\sigma_\alpha^2\mathbf 1\mathbf 1',
\end{aligned}
\label{eq:df-cov}
\end{equation}
the covariance of the transformed series $\mathbf R(\mathbf F_a\cdot)$ and $\mathbf R(\mathbf F_b\cdot)$ applied to the residual. The residual itself has $\boldsymbol\Gamma_\varepsilon=\mathcal C(\mathbf M_0,\mathbf M_0)$ and $\boldsymbol\mu_h=\operatorname{diag}(\boldsymbol\Gamma_\varepsilon)=\mathbb E[\boldsymbol\varepsilon_i\circ\boldsymbol\varepsilon_i]$.

\subsection{Random-effect weights}\label{sec:expected-theta}
The variance components of Equations~\ref{eq:re-two-way-evar}--\ref{eq:re-two-way-vvar} are quadratic forms in $\tilde{\boldsymbol\varepsilon}$, so their expectations follow from $\mathbb E[\tilde{\boldsymbol\varepsilon}'\mathbf A\tilde{\boldsymbol\varepsilon}]=\operatorname{tr}(\mathbf A\operatorname{Var}\tilde{\boldsymbol\varepsilon})$. With $\mathbf Q_T=\mathbf I-\mathbf P$ and $K_s$ slope regressors, the two-way demeaning removes $\boldsymbol\Lambda$ and the group means remove it from the group component, giving
\begin{align}
\mathbb E[\widehat\sigma_\epsilon^2]&=\frac{(N-1)\operatorname{tr}(\mathbf Q_T\boldsymbol\Omega\mathbf Q_T)}{NT\,\big(1-1/T-1/N+(1-K_s)/NT\big)},\\
\mathbb E[\widehat\sigma_{v_g}^2]&=\frac{N-1}{N-1-K_s}\,\frac{\mathbf 1'\boldsymbol\Omega\mathbf 1}{T^2}-\frac{\mathbb E[\widehat\sigma_\epsilon^2]}T,\\
\mathbb E[\widehat\sigma_{v_t}^2]&=\frac{\operatorname{tr}\!\big(\mathbf Q_T(\boldsymbol\Lambda+\boldsymbol\Omega/N)\mathbf Q_T\big)}{T-1-K_s}-\frac{\mathbb E[\widehat\sigma_\epsilon^2]}N.
\end{align}
The weights $\theta_g,\theta_t$ are Equation~\ref{eq:re-two-way-theta} evaluated at these expectations, and are then held fixed. For the simulation design below this gives $\theta_g=0.627$, $\theta_t=0.679$, against averages of $0.629$ and $0.673$ computed by paneltime at the true parameters on the simulated panels. Their sampling variation and the dependence of $\theta$ on the parameters are neglected.

\subsection{Information identity}
For one observation write $e=\varepsilon_{it}$, $v=\sigma_{it}^2$, $\mathbf d=\partial e/\partial\boldsymbol\vartheta$ and $\mathbf r=\partial v/\partial\boldsymbol\vartheta$. If $e=\sqrt v\,z$ with $z$ standard normal and independent of the information $\mathcal F$ that determines $v$, $\mathbf d$, $\mathbf r$ and the second derivatives, then the expectation of the Gaussian Hessian of Section~\ref{the-hessian} conditional on $\mathcal F$ is
\begin{equation}
\mathbb E[-\ell_{ab}\mid\mathcal F]=\frac{d_a d_b}{v}+\frac{r_a r_b}{2v^2},
\label{eq:conditional-information}
\end{equation}
since all terms that are multiplied by $e$ or by $e^2-v$ have conditional mean zero. Two exact consequences follow from the symmetry of $z$. Changing the sign of a single standardized innovation $z_s$ changes the sign of $\varepsilon_s$, but leaves every variance, every squared residual, every other residual and every derivative $d_{\cdot,s'}$ with $s'\le s$ unchanged, since these depend on squares or on earlier innovations only. Hence:
\begin{enumerate}[nosep]
\item \emph{The mean--variance block vanishes.} With $\mathbf F=\sum_m\psi_m\mathbf A_m$ the map from squared residuals to variances, $r_{a,t}=2\sum_s F_{ts}\varepsilon_s d_{a,s}$ for a mean parameter $a$, while $r_{\zeta,t}$ is unchanged by the sign of $z_s$ for a variance parameter $\zeta$. Each term of $r_ar_\zeta/v^2$ is therefore odd in one $z_s$, and its expectation is zero.
\item \emph{Only equal-time products remain in the mean block.} In
\[
r_ar_b=4\sum_{s,s'}F_{ts}F_{ts'}\varepsilon_s\varepsilon_{s'}d_{a,s}d_{b,s'}
\]
a term with $s<s'$ is odd in $z_{s'}$, and has expectation zero.
\end{enumerate}
Here $\mathbf A_m=\mathbf{GA}_\sigma^{-1}\mathbf L^m$ (Equation~\ref{eq:24}). The transform of Equation~\ref{eq:df-transform} uses group and time means, so $\varepsilon_{it}$ depends slightly on later innovations of other groups and dates. The conditional-innovation property then holds only up to terms of order $1/N$ and $1/T$, and the averaged Hessians confirm that these are negligible here: their normalized mean--variance entries are at most $0.012$ in absolute value.

\subsection{Expected variance and the ratio approximation}\label{sec:ratio-approx}
With $\mathbf g=(\sigma^2_{\mathrm{init}},\omega,\dots,\omega)'$ and $\mathbf h_i=\boldsymbol\varepsilon_i\circ\boldsymbol\varepsilon_i$, the variance is $\mathbf v_i=\mathbf{GA}_\sigma^{-1}\mathbf g+\mathbf F\mathbf h_i$ and
\begin{equation}
\bar{\mathbf v}=\mathbb E[\mathbf v_i]=\mathbf{GA}_\sigma^{-1}\mathbf g+\mathbf F\boldsymbol\mu_h.
\end{equation}
The single approximation of the calculation is
\begin{equation}
\mathbb E\!\left[\frac{x_t}{v_t}\right]\approx\frac{\mathbb E[x_t]}{\bar v_t},\qquad
\mathbb E\!\left[\frac{x_t}{v_t^2}\right]\approx\frac{\mathbb E[x_t]}{\bar v_t^2},\qquad
\mathbb E\!\left[\varepsilon_s^2 x_s\right]\approx\mu_{h,s}\,\mathbb E[x_s],
\label{eq:ratio-approximation}
\end{equation}
for the derivative products $x$. It is exact when the variance is deterministic ($\psi=0$). Otherwise it errs in both directions: $\mathbb E[1/v]>1/\mathbb E[v]$ by Jensen's inequality, but the numerators are positively correlated with $v$. The conditional-variance coefficient of variation $\sqrt{\operatorname{Var}(\sigma^2_{it})}/\sigma_\epsilon^2$ indicates its size. Evaluating the ratios exactly would require the stationary joint distribution of $(v_t,\mathbf r_t)$, which has no closed form.

\subsection{The information matrix}
\paragraph{Mean block.} The residual derivatives are linear maps of the residual: $\partial\tilde{\boldsymbol\varepsilon}_i/\partial\rho_j=-\mathbf{MA}^{-1}\mathbf L^j(\cdot)$ and $\partial\tilde{\boldsymbol\varepsilon}_i/\partial\lambda_j=-\mathbf{MA}^{-1}\mathbf L^j\mathbf M_0(\cdot)$ applied to $\mathbf u_i+\alpha_i\mathbf 1+\boldsymbol\lambda$, and $\partial\tilde{\boldsymbol\varepsilon}_i/\partial\beta_l=-\mathbf M_0\mathbf x_{i,l}$. Each is then transformed by Equation~\ref{eq:df-transform}. Writing $\mathbf D_a$ for the map of parameter $a$, the equal-time second moments are
\begin{equation}
\begin{aligned}
\mathbf p_{ab}&=\operatorname{diag}\,\mathcal C(\mathbf D_a,\mathbf D_b),\\
\mathbf p_{ll}&=\operatorname{diag}\big(\mathbf R_{\mathrm{own}}\mathbf M_0\boldsymbol\Gamma_x\mathbf M_0'\mathbf R_{\mathrm{own}}'+(N-1)\mathbf R_{\mathrm{oth}}\mathbf M_0\boldsymbol\Gamma_x\mathbf M_0'\mathbf R_{\mathrm{oth}}'\big),\\
\mathbf p_{00}&=(\mathbf S\mathbf M_0\mathbf 1)^{\circ2}
\end{aligned}
\end{equation}
for ARMA parameters, slopes and the intercept respectively (independent regressors give zero slope cross moments). Equation~\ref{eq:conditional-information} with Equation~\ref{eq:ratio-approximation} gives
\begin{equation}
\mathcal I_{ab}=N\sum_t\left[\frac{p_{ab,t}}{\bar v_t}+\frac{2}{\bar v_t^2}\sum_s F_{ts}^2\,\mu_{h,s}\,p_{ab,s}\right].
\label{eq:mean-block}
\end{equation}
The first term is the information in the residual itself. The second is the information in the future variances, through the squared residuals.

\paragraph{Variance block.} From Equation~\ref{eq:24}, $\partial\mathbf v_i/\partial\omega=\mathbf{GA}_\sigma^{-1}(0,1,\dots,1)'$ (the initial variance is fixed), $\partial\mathbf v_i/\partial\psi_m=\mathbf A_m\mathbf h_i$ and $\partial\mathbf v_i/\partial\gamma_k=\mathbf B_k\mathbf v_i$ with $\mathbf B_k=\mathbf{GA}_\sigma^{-1}\mathbf L^k$. These are affine in $\mathbf h_i$, so $\mathbb E[\mathbf r_\zeta\mathbf r_\eta']$ requires only $\boldsymbol\mu_h$ and $\mathbf Q=\mathbb E[\mathbf h_i\mathbf h_i']$, through $\mathbb E[\mathbf h\mathbf v']=\boldsymbol\mu_h\mathbf g_v'+\mathbf Q\mathbf F'$ and $\mathbb E[\mathbf v\mathbf v']=\mathbf g_v\bar{\mathbf v}'+\mathbf F\boldsymbol\mu_h\mathbf g_v'+\mathbf F\mathbf Q\mathbf F'$, where $\mathbf g_v=\mathbf{GA}_\sigma^{-1}\mathbf g$. Then
\begin{equation}
\mathcal I_{\zeta\eta}=\frac N2\sum_t\frac{\mathbb E[r_{\zeta,t}r_{\eta,t}]}{\bar v_t^2},
\qquad \mathcal I_{a\zeta}=0.
\label{eq:garch-information}
\end{equation}

\paragraph{Fourth moments.} For a Gaussian vector, Isserlis' theorem gives $\mathbb E[\varepsilon_s^2\varepsilon_t^2]=\Gamma_{\varepsilon,ss}\Gamma_{\varepsilon,tt}+2\Gamma_{\varepsilon,st}^2$. GARCH innovations are not Gaussian, and their excess kurtosis is added as follows. Following \citet{bollerslev1986generalized}, $\epsilon_{it}^2-\sigma_\epsilon^2$ is an ARMA process with autoregressive coefficients $\psi_j+\gamma_j$ and moving-average coefficients $-\gamma_j$, driven by $\eta_{it}=\epsilon_{it}^2-\sigma_{it}^2$. With $S=\sum_{j\ge0}\pi_j^2$ the sum of its squared MA($\infty$) weights,
\begin{equation}
\begin{gathered}
\operatorname{Var}(\eta_{it})=\frac{2\sigma_\epsilon^4}{3-2S},\qquad
\operatorname{Var}(\sigma_{it}^2)=(S-1)\operatorname{Var}(\eta_{it}),\\
\mathbb{E}[\epsilon_{it}^4]<\infty\iff S<\tfrac32,
\end{gathered}
\label{eq:fourth-moment-condition}
\end{equation}
which for GARCH(1,1) is $3\psi^2+2\psi\gamma+\gamma^2<1$. Let $\Delta(k)=\operatorname{Cov}(\epsilon_t^2,\epsilon_{t-k}^2)-2\sigma_\epsilon^4\mathbb 1[k=0]$ be the excess over Gaussian white noise, and $\boldsymbol\Delta=[\Delta(|r-r'|)]$. The residual is a sum of independent components $\mathbf C\boldsymbol\epsilon$, namely $\mathbf C=\mathbf R_{\mathrm{own}}\mathbf M_0\mathbf W$ for the own group and $\mathbf C=\mathbf R_{\mathrm{oth}}\mathbf M_0\mathbf W$ for each of the $N-1$ other groups, where $\mathbf W$ is the MA($\infty$) matrix of $u$, plus Gaussian effects. Expanding $\mathbb E[\varepsilon_s^2\varepsilon_t^2]$ for symmetric innovations gives
\begin{equation}
\begin{aligned}
\mathbf Q={}&\boldsymbol\mu_h\boldsymbol\mu_h'+2\boldsymbol\Gamma_\varepsilon^{\circ2}
+\sum_{\mathbf C}\Big[(\mathbf C\circ\mathbf C)\boldsymbol\Delta(\mathbf C\circ\mathbf C)'-2\Delta(0)(\mathbf C\circ\mathbf C)(\mathbf C\circ\mathbf C)'\\
&+2\big[(\mathbf c_s\circ\mathbf c_t)'\boldsymbol\Delta(\mathbf c_s\circ\mathbf c_t)\big]_{s,t}\Big],
\end{aligned}
\label{eq:fourth-moment-excess}
\end{equation}
with $\mathbf c_s$ row $s$ of $\mathbf C$.

\paragraph{Intercept.} In the design, the sample means $\bar\alpha$ and $\bar\lambda$ of the effects are not separable from $\beta_0$. paneltime therefore adds the estimated $\sigma_\alpha^2/N+\sigma_\lambda^2/T$ to the variance of the intercept. The data-free counterpart adds the true values:
\begin{equation}
\operatorname{SE}_{\beta_0}=\sqrt{\big[\mathcal I^{-1}\big]_{00}+\sigma_\alpha^2/N+\sigma_\lambda^2/T}.
\end{equation}

\subsection{Accuracy in the simulation design}\label{sec:expected-se-truth}
Table~\ref{tab:data-free-se} compares the data-free SEs with three simulated quantities for the simulation design ($N=20$, $T=100$, ARMA(2,2)--GARCH(2,2), unit effect variances, 100 panels). These are diagnostics only: $\mathbb E[-\mathbf H]$ is paneltime's negative Hessian at the true parameters averaged over the simulated panels, so it is the exact counterpart of Equation~\ref{eq:fisher-definition} up to simulation noise. SD is the standard deviation of the estimates. Median SE is the median robust standard error reported by paneltime.

\begin{table}[H]
\centering
\caption{Data-free SEs against simulated references}\label{tab:data-free-se}
\begin{tabular}{lrrrr}
\toprule
Parameter & Data-free & $\mathbb E[-\mathbf H]$ & SD & Median SE\\
\midrule
Intercept$^{a}$ & 0.262 & 0.262 & 0.254 & 0.288\\
$X_0$ & 0.045 & 0.043 & 0.051 & 0.045\\
$X_1$ & 0.045 & 0.043 & 0.045 & 0.045\\
$\rho_1$ & 0.136 & 0.144 & 0.201 & 0.148\\
$\rho_2$ & 0.070 & 0.075 & 0.102 & 0.071\\
$\lambda_1$ & 0.136 & 0.144 & 0.202 & 0.149\\
$\lambda_2$ & 0.054 & 0.057 & 0.072 & 0.066\\
$\gamma_1$ & 0.205 & 0.175 & 0.170 & 0.177\\
$\gamma_2$ & 0.177 & 0.152 & 0.132 & 0.164\\
$\psi_1$ & 0.021 & 0.025 & 0.023 & 0.028\\
$\psi_2$ & 0.026 & 0.029 & 0.025 & 0.031\\
$\omega$ & 0.120 & 0.120 & 0.186 & 0.151\\
\bottomrule
\end{tabular}

\smallskip
{\footnotesize $^{a}$ Both columns include $\sigma_\alpha^2/N+\sigma_\lambda^2/T=0.06$.}
\end{table}

The mean block and $\omega$ are within about $6\%$ of $\mathbb E[-\mathbf H]$. The GARCH and ARCH entries differ by up to $17\%$ in opposite directions. The variance-block differences are mainly due to the ratio approximation of Equation~\ref{eq:ratio-approximation}, the only error that does not vanish with the sample size. A first-order (delta) correction of the ratios using $\mathbf Q$ was tested. It improved individual information entries, but reduced the average absolute SE error only from $9.4\%$ to $8.7\%$, and is not used.

The larger gaps between $\mathbb E[-\mathbf H]$ and the sampling SD for the ARMA parameters and $\omega$ are not errors of the data-free calculation. They reflect the non-quadratic likelihood at $T=100$, and no expected-information calculation can capture them. For $\omega$ the reported SEs are right-skewed (mean $0.192$, median $0.151$) because they are evaluated at estimates with low persistence, where the likelihood is flatter. Their information average $\big(\overline{\operatorname{SE}^{-2}}\big)^{-1/2}=0.130$ is close to the data-free value.

\paragraph{Special case.} With $\psi=\gamma=0$ the variance is deterministic and Equation~\ref{eq:ratio-approximation} is exact. Without effects or ARMA terms the calculation then reduces to the familiar $\operatorname{SE}(\widehat\beta_j)=\sqrt{\omega/(NT\sigma_x^2)}$ and $\operatorname{SE}(\widehat\omega)=\omega\sqrt{2/[N(T-1)]}$, where $T-1$ reflects the fixed initial variance.

\subsection{Time-series specialization}\label{sec:df-time}
For the time-series design, $N=1$ and $\sigma_\alpha^2=\sigma_\lambda^2=0$, so the effect transform disappears. The regressors and disturbance have the same ARMA law but independent innovations, with variances $\sigma_x^2$ and $\sigma_\epsilon^2$. Write
\[
A(z)=1-\sum_j\rho_jz^j,\qquad B(z)=1+\sum_j\lambda_jz^j,\qquad
G(z)=1-\sum_j\gamma_jz^j.
\]
Stationary residual derivatives for the AR and MA parameters have filters
$D_{\rho_j}(z)=-z^j/A(z)$ and $D_{\lambda_j}(z)=-z^j/B(z)$ applied to $\epsilon_t$.
Consequently, with $D_{a,h}$ denoting lag coefficients,
\begin{equation}
p_{ab}=\sigma_\epsilon^2\sum_{h\ge0}D_{a,h}D_{b,h},\qquad
p_{00}=\left(\frac{A(1)}{B(1)}\right)^2,\qquad p_{ll}=\sigma_x^2.
\label{eq:df-time-mean-moments}
\end{equation}
All remaining mean-parameter cross moments vanish for the independent, zero-mean regressor design.

It is important to retain the finite GARCH initialization: replacing it by a fully stationary variance-derivative process can substantially change the SEs in this design. Let $b_h$ and $f_h$ be the coefficients of
$b(z)=1/G(z)$ and $f(z)=b(z)\sum_j\psi_jz^j$, and index the observations by $t=0,\ldots,T-1$. With $g_0=\sigma_\epsilon^2$ and $g_t=\omega$ for $t>0$,
\begin{equation}
\bar v_t=\sum_{h=0}^t b_hg_{t-h}+\sigma_\epsilon^2\sum_{h=0}^t f_h,\qquad
\mathcal I_{ab}\approx p_{ab}\sum_{t=0}^{T-1}
\left[\frac1{\bar v_t}+
\frac{2\sigma_\epsilon^2}{\bar v_t^2}\sum_{h=0}^t f_h^2\right].
\label{eq:df-time-mean-information}
\end{equation}
This is Equation~\ref{eq:mean-block} without the panel transform.

The variance derivatives are finite convolutions of $h_t=\epsilon_t^2$. Their random-part filters and means are
\begin{equation}
\begin{aligned}
c_{\psi_j}(z)&=z^jb(z),&
M_{\psi_j,t}&=\sigma_\epsilon^2\sum_{h=0}^t c_{\psi_j,h},\\
c_{\gamma_j}(z)&=z^jb(z)f(z),&
M_{\gamma_j,t}&=\sum_{h=0}^t [z^jb(z)]_h\,\bar v_{t-h},\\
c_\omega(z)&=0,&
M_{\omega,t}&=\sum_{h=0}^{t-1}b_h .
\end{aligned}
\label{eq:df-time-derivatives}
\end{equation}
The empty sum at $t=0$ is zero, since the initial variance is fixed. If
$H(k)=\operatorname{Cov}(\epsilon_t^2,\epsilon_{t-k}^2)$ is obtained from the squared-innovation ARMA representation above, then
\begin{equation}
\begin{aligned}
V_{ab,t}&=\sum_{i=0}^t\sum_{j=0}^t c_{a,i}c_{b,j}H(|i-j|),\\
\mathcal I_{ab}&\approx\frac12\sum_{t=0}^{T-1}
\frac{M_{a,t}M_{b,t}+V_{ab,t}}{\bar v_t^2}.
\end{aligned}
\label{eq:df-time-variance-information}
\end{equation}
These sums can be evaluated by lag convolutions and cumulative sums, without panel matrices or simulated observations.
The implementation uses stationary innovation and ARMA-derivative moments, omitting the transient in residuals caused by the zero-presample ARMA filter. It retains the fixed initial variance and finite GARCH derivative paths. Thus, for pure GARCH the moments coincide with the generic finite-initialization calculation; with ARMA there is an additional boundary approximation, besides the GARCH ratio/product approximation of Equation~\ref{eq:ratio-approximation}. These are expected-information SEs, not exact finite-sample or robust SEs.

\subsection{Static balanced panels}\label{sec:df-static-panel}
With no ARMA or GARCH terms, put $\sigma^2=\sigma_\epsilon^2=\omega$ and let
$d_w=(N-1)(T-1)$. The grand-mean-preserving effect transform has eigenvalues
$1$, $1-\theta_g$, $1-\theta_t$ and $1$ on the grand-mean, group-contrast,
time-contrast and within-panel subspaces, whose dimensions are respectively
$1$, $N-1$, $T-1$ and $d_w$. Therefore
\begin{equation}
d_R=\operatorname{tr}(R'R)
=d_w+(N-1)(1-\theta_g)^2+(T-1)(1-\theta_t)^2+1.
\label{eq:df-static-trace}
\end{equation}
For independent zero-mean regressors of variance $\sigma_x^2$, the closed-form information and SE references are
\begin{equation}
\begin{aligned}
\mathcal I_{\beta_l\beta_l}&=\frac{\sigma_x^2d_R}{\sigma^2},&
\operatorname{SE}_{\rm DF}(\widehat\beta_l)&=\sqrt{\frac{\sigma^2}{\sigma_x^2d_R}},\\
\mathcal I_{\beta_0\beta_0}&=\frac{NT}{\sigma^2},&
\mathcal I_{\omega\omega}&=\frac{N(T-1)}{2\sigma^4}.
\end{aligned}
\label{eq:df-static-information}
\end{equation}
For two-way fixed effects $\theta_g=\theta_t=1$, so $d_R=d_w+1$.
The extra $1$ is the grand-mean direction retained by paneltime's transform; it disappears relative to $d_w$ asymptotically. The intercept reference
$\sigma/\sqrt{NT}$ is conditional on centered fixed effects. It is not the across-panel sampling SD when new random effects are generated in every replication.

For random effects, the static versions of the expected variance components are
\begin{equation}
\begin{aligned}
E_e&=\frac{\sigma^2d_w}{d_w-K_s},\\
E_g&=\frac{N-1}{N-1-K_s}\left(\sigma_\alpha^2+\frac{\sigma^2}{T}\right)-\frac{E_e}{T},\\
E_t&=\frac{T-1}{T-1-K_s}\left(\sigma_\lambda^2+\frac{\sigma^2}{N}\right)-\frac{E_e}{N}.
\end{aligned}
\label{eq:df-static-components}
\end{equation}
Set negative $E_g,E_t$ to zero, evaluate Equation~\ref{eq:re-two-way-theta} at these values, and hold the resulting weights fixed. The random-effects intercept reference is
\begin{equation}
\operatorname{SE}_{\rm DF}(\widehat\beta_0)=
\sqrt{\frac{\sigma^2}{NT}+\frac{\sigma_\alpha^2}{N}+\frac{\sigma_\lambda^2}{T}}.
\label{eq:df-static-intercept}
\end{equation}
There is no stochastic-variance ratio approximation in these static formulas. They are exact for the expected quadratic information with the specified fixed transform weights and true variance, not exact finite-sample sampling variances: inversion after averaging the design information differs from averaging its inverse, and estimation and differentiation of random-effect weights remain outside this reference.
For the panel design ($N=20$, $T=100$, $K_s=2$, $\sigma=2$, $\sigma_x=1$, unit effect variances), the slope SEs are $0.04610$ for fixed effects and $0.04590$ for random effects; the intercept references are $0.04472$ conditional on fixed effects and $0.24900$ with random effects.

\section{Starting values}\label{sec:starting-values}

The joint likelihood is optimized from one deterministic starting vector by
default. Initialization is not a separate estimate reported to the user, and
does not alter the derivatives, Hessian, or parameter restrictions of the final
likelihood. Explicit user-supplied starting coefficients take precedence.

\subsection{Preliminary regression and full-order ARMA fit}

First, ordinary least squares is applied to the differenced data, after the
existing preliminary group/time effects transformations. Write these working
variables as $\widetilde y_{it}$ and $\widetilde{\mathbf x}_{it}$ and put
$u_{it}=\widetilde y_{it}-\widetilde{\mathbf x}_{it}'\boldsymbol\beta$.
The effects transformation is held fixed during this preliminary fit. In
particular, this is an initialization approximation, not the final likelihood:
the final estimator removes effects \emph{after} ARMA filtering and recomputes
random-effect weights as described in Section~\ref{sec:effects-implementation}.
With $d>0$, the first $d$ differencing-boundary observations are excluded from
the preliminary fit.

For requested orders $(p,q)$, a Hannan--Rissanen-style two-regression start
\cite{hannan1982recursive} supplies all AR and MA coefficients. An auxiliary
AR($r$) regression of the OLS residuals gives approximate innovations
$\widehat a_{it}$. A second pooled regression estimates
\[
 u_{it}=\sum_{j=1}^{p}\rho_j u_{i,t-j}
       +\sum_{j=1}^{q}\lambda_j\widehat a_{i,t-j}
       +\text{error}_{it}.
\]
Only complete lag windows within an individual are used; no lag crosses a
group boundary or an excluded observation. The auxiliary order is
$r=\min\{20,\max[2\max(p,q)+1,\lceil\log T_{\max}\rceil]\}$, reduced when
necessary to leave more usable rows than regression columns. If the auxiliary
regressions cannot be identified, a warning is issued and zero ARMA
coefficients are used to seed the subsequent conditional-sum-of-squares fit.

The regression and all requested ARMA coefficients are then jointly refined by
minimizing
\begin{align}
 e_{it}
 &=u_{it}-\sum_{j=1}^{p}\rho_j u_{i,t-j}
          -\sum_{j=1}^{q}\lambda_j e_{i,t-j},\\
 Q(\boldsymbol\beta,\boldsymbol\rho,\boldsymbol\lambda)
 &=\frac{1}{n_c}\sum_{(i,t)\in\mathcal C}e_{it}^{\,2}.
\end{align}
Here presample residuals and innovations are zero, the recursion restarts at
each group or excluded-observation boundary, and $\mathcal C$ omits the first
$p$ observations of each contiguous run, as in conditional least squares.
GARCH is deliberately ignored at this stage.
The implementation evaluates the inverse-MA filtering of innovations and
their derivative channels through the existing native ARMA recursion,
packing contiguous runs as independent series. This changes only the
computation, not the criterion or its boundary conditions.

Stationarity and invertibility of the preliminary ARMA coefficients are
enforced through reflection coefficients \cite{monahan1984note}.
For an unconstrained vector $\boldsymbol\eta$, let $\kappa_j=\tanh\eta_j$ and
construct
\[
 a_j^{(j)}=\kappa_j,\qquad
 a_\ell^{(j)}=a_\ell^{(j-1)}
             -\kappa_j a_{j-\ell}^{(j-1)},\quad 1\leq\ell<j.
\]
Use $\boldsymbol\rho=\mathbf a^{(p)}$ and
$\boldsymbol\lambda=-\mathbf a^{(q)}$, reflecting the MA sign convention
$1+\sum_j\lambda_j L^j$. Auxiliary roots outside a radius of $0.98$ are
reflected inside that radius before conversion to reflection coordinates.
These are safeguards on the \emph{preliminary} fit only, not new restrictions
on the final likelihood.

The CSS refinement uses damped Gauss--Newton steps, with at most 10 iterations,
and analytic first derivatives. It deliberately stops well before a precise
standalone ARMA estimate is needed: in internally rescaled coordinates it
stops when $\|J'\mathbf e/n_c\|_\infty\leq 10^{-4}$ or an accepted step reduces
$Q$ by at most $10^{-4}\max(Q,10^{-12})$, where $J$ is the innovation Jacobian.
The joint likelihood optimizer subsequently refines these starting values.
For example, before the reflection-coordinate
chain rule,
\[
 \frac{\partial e_{it}}{\partial\boldsymbol\beta}
 =-\widetilde{\mathbf x}_{it}
   +\sum_{j=1}^p\rho_j\widetilde{\mathbf x}_{i,t-j}
   -\sum_{j=1}^q\lambda_j
       \frac{\partial e_{i,t-j}}{\partial\boldsymbol\beta},
\]
and the direct AR and MA derivative inputs are respectively $-u_{i,t-j}$
and $-e_{i,t-j}$, followed by the same inverse-MA recursion.
Variables are rescaled internally to improve conditioning and converted back
afterwards. Only objective-decreasing steps are accepted. Damping increases
when the actual objective reduction is less than one quarter of the
Gauss--Newton prediction, decreases when it exceeds three quarters, and is
otherwise retained. This avoids repeatedly accepting overly aggressive steps
that make little progress. An iteration-limit
or numerical-stall warning distinguishes a bounded preliminary search from
a converged fit; the best accepted preliminary coefficients remain usable as
starting values. These initialization-specific warnings are reported without
aborting the final likelihood fit; unrelated numerical warnings remain errors.

\subsection{Variance starts and relation to rugarch}

Let $s_e^2$ be the sample variance of the full-order preliminary innovations
over included observations, including the AR conditioning boundary. The
existing squared-residual moment rule supplies the first ARCH and GARCH
starts, $\psi_1^{(0)}$ and $\gamma_1^{(0)}$; higher variance lags initially
remain zero. For the ordinary GARCH specification,
\[
 v_{\mathrm{initial}}=s_e^2,\qquad
 \omega^{(0)}
 =h(0,s_e^2,s_e^2)\bigl(1-\psi_1^{(0)}-\gamma_1^{(0)}\bigr).
\]
Thus the initial innovation variance no longer relies on an ARMA(1,1)
variance formula when higher orders were requested. EGARCH retains its
existing transformed-variance initialization. With no ARMA terms, the
existing OLS and variance initialization are unchanged.

The broad strategy is shared with \texttt{rugarch}
\cite{ghalanos2024rugarch}: its version 1.5-5 initializes the mean equation by
calling \texttt{stats::arima(method="CSS")} at the full requested order
\cite{rcore2026arima}. That routine starts ARMA coefficients at zero and uses
BFGS, without its stationarity transformation on the CSS path. Paneltime
instead uses the auxiliary-regression seed, constrained reflection coordinates,
and a short damped least-squares refinement, extended to pooled panel runs.
It does not copy rugarch's GARCH starts: for GARCH(2,2), rugarch defaults to
ARCH coefficients $(0.025,0.025)$, GARCH coefficients $(0.45,0.45)$, and
a variance constant equal to the variance of its input series divided by
1000 (after scaling, if requested).
Neither preliminary procedure guarantees the global maximum of the joint
ARMA--GARCH likelihood.

\section{Extensions}

\subsection{Tobit model}\label{tobit-model}

For censored observations, i.e.\ where the underlying, unobserved dependent variable lies beyond a threshold \citep{tobin1958estimation,amemiya1973regression}, the contribution to the likelihood is given by the normal cumulative distribution function evaluated at the threshold,

\begin{equation}
F\left( \gamma\frac{\epsilon_{i,t}}{\sigma_{i,t}} \right) = \int_{- \infty}^{\gamma\epsilon_{i,t}/\sigma_{i,t}}{\frac{1}{\sqrt{2\pi}}\exp\left( - \frac{1}{2}\varepsilon^{2} \right)d\varepsilon}
\label{eq:censored}
\end{equation}

where $\epsilon_{i,t}$ is the residual evaluated at the censoring threshold. Let
$\gamma$ indicate the type of threshold, with $\gamma = 1$ for a
lower threshold and $\gamma = -1$ for an upper threshold. Define $x=\gamma\epsilon_{i,t}/\sigma_{i,t}$, $\lambda(x)=f(x)/F(x)$, and $s=\ln\sigma_{i,t}$. The derivatives with respect to residual and log standard deviation are

\begin{equation}
\frac{\partial\ln F(x)}{\partial\epsilon_{i,t}} = \frac{\gamma\lambda(x)}{\sigma_{i,t}}
\label{eq:tobit_deriv1}
\end{equation}

\begin{equation}
\frac{\partial\ln F(x)}{\partial s}=-x\lambda(x)
\label{eq:tobit_deriv2}
\end{equation}

\begin{equation}
\frac{\partial^{2}\ln F(x)}{\partial\epsilon_{i,t}^{2}}=-\left[\gamma\lambda(x)\frac{\epsilon_{i,t}}{\sigma_{i,t}^{3}}+\frac{\lambda(x)^{2}}{\sigma_{i,t}^{2}}\right]
\label{eq:tobit_deriv3}
\end{equation}

\begin{equation}
\frac{\partial^{2}\ln F(x)}{\partial\epsilon_{i,t}\,\partial s}
=\frac{\gamma\lambda(x)}{\sigma_{i,t}}\left[x^{2}+x\lambda(x)-1\right]
\label{eq:tobit_deriv4}
\end{equation}

\begin{equation}
\frac{\partial^{2}\ln F(x)}{\partial s^{2}}
=x\lambda(x)-x^{3}\lambda(x)-x^{2}\lambda(x)^{2}
\label{eq:tobit_deriv5}
\end{equation}

The optimizer's log-variance coordinate is $r=\ln\sigma_{i,t}^{2}=2s$. Consequently,
\begin{equation}
\frac{\partial L}{\partial r}=\frac{1}{2}\frac{\partial L}{\partial s},\qquad
\frac{\partial^{2}L}{\partial\epsilon\,\partial r}=\frac{1}{2}\frac{\partial^{2}L}{\partial\epsilon\,\partial s},\qquad
\frac{\partial^{2}L}{\partial r^{2}}=\frac{1}{4}\frac{\partial^{2}L}{\partial s^{2}}.
\label{eq:tobit-logvar-conversion}
\end{equation}
When the variance state is $v=\sigma^{2}$ rather than $r$, the conversion in Equation~\ref{eq:variance-logvariance-chain} is also required. The current implementation does not consistently perform this conversion for Tobit observations in ordinary GARCH, and its coded $\partial^{2}L/\partial r^{2}$ term has the wrong sign on the $x^{2}\lambda(x)^{2}$ component. Thus the formulas above are the analytic reference; they do not describe the current Tobit Hessian implementation exactly.



\subsection{Instrumental variables}\label{instrumental-variables}

The current instrumental-variable code path does not implement the conventional two-stage least-squares projection described by $\mathbf{P}_{Z}\mathbf{X}$. When a non-constant instrument matrix is supplied, \texttt{Panel.set\_instrumentals()} assigns that matrix directly as \texttt{XIV}, and the likelihood forms its regression residual using \texttt{XIV}. Thus the current implementation uses the supplied instrument columns in place of the regressor design for the likelihood transformation; it does not estimate first-stage fitted regressors. The earlier description as standard 2SLS was inaccurate. This behavior should not be interpreted as a general IV estimator.

\bibliography{literature}
\end{document}